% =====================================================================
%  libreta.tex
%  ========================
%
%
%  Descripción:
%  ------------
%  Apuntes de clase del curso de Aprendizaje Computacional
%
%
%  Metadata:
%  ----------
%   * Author : zxxz6 (Bryan Violante Arriaga)
%   * Version: 1.6.0
%
%
%  History:
%  ------------
%  Author      Date            Description
%  zxxz6       20/09/2026      Clase 17 Sep, gradiente, delta y backprop
%  zxxz6       20/09/2026      Clase 15 Sep, TLU y separabilidad lineal
%  zxxz6       14/09/2026      Clase 10 Sep, preparacion de datos y redes
%  zxxz6       09/09/2026      Clase 8 Sep, forma dual, kernels y multiclase
%  zxxz6       09/09/2026      Complete contra las diapositivas 03 a 05
%  zxxz6       03/09/2026      Pase a LaTeX la toma cruda del 03/09
%  zxxz6       01/09/2026      Pase a LaTeX la toma cruda del 01/09
%  zxxz6       27/08/2026      Pase a LaTeX la toma cruda del 27/08
%  zxxz6       25/08/2026      Creation
%
%
% =====================================================================
\documentclass[11pt,letterpaper]{article}

% =====================================================================
%  DATOS DE LA MATERIA
% =====================================================================
\newcommand{\materia}{Aprendizaje Computacional}
\newcommand{\profesor}{Dr. José Francisco Martínez Trinidad}
\newcommand{\periodo}{Otoño 2026}
\newcommand{\alumno}{Bryan Ezequiel Violante Arriaga}

% =====================================================================
%  Preambulo.tex
%  ========================
%
%
%  Descripción:
%  ------------
%  Preámbulo compartido de los apuntes de clase: paquetes, estilo de
%  página, entornos de teorema, las cuatro cajas de color, el estilo de
%  listings y el pseudocódigo en español. Se carga con \input desde
%  cada documento en lugar de copiarlo, para que un arreglo aquí llegue
%  a todas las libretas de una vez.
%
%
%  Considerations:
%  ------------
%  - Se carga con \input y ruta relativa, no con \usepackage, porque
%    \usepackage solo busca en la carpeta del documento y en el árbol
%    de TeX. Un .sty instalado en TEXMFHOME evitaría la ruta, pero vive
%    fuera del repo y una clona nueva no compilaría
%  - Los cuatro datos de la materia se declaran en el documento ANTES
%    del \input. hyperref expande \materia al cargarse, así que si no
%    existe todavía la compilación truena
%  - Cuando la materia la imparte más de un profesor, los nombres van
%    en un solo \profesor separados por \\. La portadilla los imprime
%    dentro de un center, que es donde \\ sí corta renglón. Repetir
%    \newcommand{\profesor} no es opción: truena con "already defined"
%  - algorithmicx no tiene opción de idioma y babel no lo traduce, así
%    que las palabras clave del pseudocódigo se renombran a mano
%  - Los nombres de los entornos van sin acento (\begin{definicion})
%    aunque el rótulo impreso sí lo lleve
%  - babel se carga con es-noshorthands porque en español el " queda
%    activo y se come el texto que le sigue, además de romper listings
%
%
%  Metadata:
%  ----------
%   * Author : zxxz6 (Bryan Violante Arriaga)
%   * Version: 1.7.0
%
%
%  History:
%  ------------
%  Author      Date            Description
%  zxxz6       24/09/2026      El nombre del alumno en el encabezado
%  zxxz6       24/09/2026      Encabezado de una hoja de un ejercicio
%  zxxz6       23/09/2026      Macros de los documentos de soluciones
%  zxxz6       30/08/2026      Macros de lim, limsup y liminf sobre n
%  zxxz6       25/08/2026      Macros de Omega, Theta, o y omega chicas
%  zxxz6       20/08/2026      Cargue pgfplots para graficas de datos
%  zxxz6       20/08/2026      Cargue tikz para diagramas de conjuntos
%  zxxz6       19/08/2026      Documente \profesor con varios nombres
%  zxxz6       19/08/2026      Agregue la caja Idea Random
%  zxxz6       18/08/2026      Creation
%
%
% =====================================================================
%  USO
%  ---
%  \documentclass[11pt,letterpaper]{article}
%  \newcommand{\materia}{...}      % los cuatro datos van arriba
%  \newcommand{\profesor}{...}
%  \newcommand{\periodo}{...}
%  \newcommand{\alumno}{...}
%  % =====================================================================
%  Preambulo.tex
%  ========================
%
%
%  Descripción:
%  ------------
%  Preámbulo compartido de los apuntes de clase: paquetes, estilo de
%  página, entornos de teorema, las cuatro cajas de color, el estilo de
%  listings y el pseudocódigo en español. Se carga con \input desde
%  cada documento en lugar de copiarlo, para que un arreglo aquí llegue
%  a todas las libretas de una vez.
%
%
%  Considerations:
%  ------------
%  - Se carga con \input y ruta relativa, no con \usepackage, porque
%    \usepackage solo busca en la carpeta del documento y en el árbol
%    de TeX. Un .sty instalado en TEXMFHOME evitaría la ruta, pero vive
%    fuera del repo y una clona nueva no compilaría
%  - Los cuatro datos de la materia se declaran en el documento ANTES
%    del \input. hyperref expande \materia al cargarse, así que si no
%    existe todavía la compilación truena
%  - Cuando la materia la imparte más de un profesor, los nombres van
%    en un solo \profesor separados por \\. La portadilla los imprime
%    dentro de un center, que es donde \\ sí corta renglón. Repetir
%    \newcommand{\profesor} no es opción: truena con "already defined"
%  - algorithmicx no tiene opción de idioma y babel no lo traduce, así
%    que las palabras clave del pseudocódigo se renombran a mano
%  - Los nombres de los entornos van sin acento (\begin{definicion})
%    aunque el rótulo impreso sí lo lleve
%  - babel se carga con es-noshorthands porque en español el " queda
%    activo y se come el texto que le sigue, además de romper listings
%
%
%  Metadata:
%  ----------
%   * Author : zxxz6 (Bryan Violante Arriaga)
%   * Version: 1.7.0
%
%
%  History:
%  ------------
%  Author      Date            Description
%  zxxz6       24/09/2026      El nombre del alumno en el encabezado
%  zxxz6       24/09/2026      Encabezado de una hoja de un ejercicio
%  zxxz6       23/09/2026      Macros de los documentos de soluciones
%  zxxz6       30/08/2026      Macros de lim, limsup y liminf sobre n
%  zxxz6       25/08/2026      Macros de Omega, Theta, o y omega chicas
%  zxxz6       20/08/2026      Cargue pgfplots para graficas de datos
%  zxxz6       20/08/2026      Cargue tikz para diagramas de conjuntos
%  zxxz6       19/08/2026      Documente \profesor con varios nombres
%  zxxz6       19/08/2026      Agregue la caja Idea Random
%  zxxz6       18/08/2026      Creation
%
%
% =====================================================================
%  USO
%  ---
%  \documentclass[11pt,letterpaper]{article}
%  \newcommand{\materia}{...}      % los cuatro datos van arriba
%  \newcommand{\profesor}{...}
%  \newcommand{\periodo}{...}
%  \newcommand{\alumno}{...}
%  % =====================================================================
%  Preambulo.tex
%  ========================
%
%
%  Descripción:
%  ------------
%  Preámbulo compartido de los apuntes de clase: paquetes, estilo de
%  página, entornos de teorema, las cuatro cajas de color, el estilo de
%  listings y el pseudocódigo en español. Se carga con \input desde
%  cada documento en lugar de copiarlo, para que un arreglo aquí llegue
%  a todas las libretas de una vez.
%
%
%  Considerations:
%  ------------
%  - Se carga con \input y ruta relativa, no con \usepackage, porque
%    \usepackage solo busca en la carpeta del documento y en el árbol
%    de TeX. Un .sty instalado en TEXMFHOME evitaría la ruta, pero vive
%    fuera del repo y una clona nueva no compilaría
%  - Los cuatro datos de la materia se declaran en el documento ANTES
%    del \input. hyperref expande \materia al cargarse, así que si no
%    existe todavía la compilación truena
%  - Cuando la materia la imparte más de un profesor, los nombres van
%    en un solo \profesor separados por \\. La portadilla los imprime
%    dentro de un center, que es donde \\ sí corta renglón. Repetir
%    \newcommand{\profesor} no es opción: truena con "already defined"
%  - algorithmicx no tiene opción de idioma y babel no lo traduce, así
%    que las palabras clave del pseudocódigo se renombran a mano
%  - Los nombres de los entornos van sin acento (\begin{definicion})
%    aunque el rótulo impreso sí lo lleve
%  - babel se carga con es-noshorthands porque en español el " queda
%    activo y se come el texto que le sigue, además de romper listings
%
%
%  Metadata:
%  ----------
%   * Author : zxxz6 (Bryan Violante Arriaga)
%   * Version: 1.7.0
%
%
%  History:
%  ------------
%  Author      Date            Description
%  zxxz6       24/09/2026      El nombre del alumno en el encabezado
%  zxxz6       24/09/2026      Encabezado de una hoja de un ejercicio
%  zxxz6       23/09/2026      Macros de los documentos de soluciones
%  zxxz6       30/08/2026      Macros de lim, limsup y liminf sobre n
%  zxxz6       25/08/2026      Macros de Omega, Theta, o y omega chicas
%  zxxz6       20/08/2026      Cargue pgfplots para graficas de datos
%  zxxz6       20/08/2026      Cargue tikz para diagramas de conjuntos
%  zxxz6       19/08/2026      Documente \profesor con varios nombres
%  zxxz6       19/08/2026      Agregue la caja Idea Random
%  zxxz6       18/08/2026      Creation
%
%
% =====================================================================
%  USO
%  ---
%  \documentclass[11pt,letterpaper]{article}
%  \newcommand{\materia}{...}      % los cuatro datos van arriba
%  \newcommand{\profesor}{...}
%  \newcommand{\periodo}{...}
%  \newcommand{\alumno}{...}
%  \input{../../templates/Preambulo.tex}
%  \begin{document}
%  \portadilla
%
%  Con dos o mas profesores se declara UN solo \profesor y los nombres
%  se separan con \\, uno por renglon:
%
%  \newcommand{\profesor}{Dra. Fulana de Tal \\
%                         Dr. Mengano de Cual}
% =====================================================================


% ==================== IDIOMA Y CODIFICACIÓN ====================
\usepackage[utf8]{inputenc}
\usepackage[T1]{fontenc}
\usepackage[spanish,es-tabla,es-noshorthands]{babel}
\usepackage{lmodern}
\usepackage{microtype}

% ==================== PÁGINA ====================
\usepackage[letterpaper,margin=2.3cm,headheight=14pt]{geometry}
\usepackage{parskip}
\usepackage{fancyhdr}
\usepackage{enumitem}

% ==================== MATEMÁTICAS ====================
\usepackage{amsmath,amssymb,amsthm,mathtools}

% ==================== FIGURAS, TABLAS, CÓDIGO ====================
\usepackage{graphicx}
\usepackage{booktabs}
\usepackage{array}
% float da el especificador [H], que clava una tabla donde se
% escribio. Una traza que LaTeX manda dos paginas adelante deja
% de leerse junto al paso que la produjo
\usepackage{float}
\usepackage{xcolor}
% tikz ya viene arrastrado por tcolorbox[most], pero se carga
% explicito para no depender de ese detalle; la libreria babel
% protege a tikz de los caracteres activos del español
\usepackage{tikz}
\usetikzlibrary{babel}
% pgfplots dibuja graficas de datos (ejes, escalas log) sobre tikz
\usepackage{pgfplots}
\pgfplotsset{compat=1.18}
\usepackage{listings}
\usepackage[ruled,section]{algorithm}
\usepackage{algpseudocode}
\usepackage[most]{tcolorbox}
% soul da \hl, el unico resaltado que parte renglon; \colorbox no
\usepackage{soul}
\usepackage{hyperref}

% =====================================================================
%  ESTILO
% =====================================================================

% --- encabezado y pie de cada página ---
\pagestyle{fancy}
\fancyhf{}
\fancyhead[L]{\small\textsc{\materia}}
\fancyhead[R]{\small\periodo}
\fancyfoot[C]{\small\thepage}
\renewcommand{\headrulewidth}{0.4pt}

% --- enlaces ---
\hypersetup{
  colorlinks=true,
  linkcolor=black,
  urlcolor=blue,
  citecolor=black,
  pdftitle={Apuntes de \materia},
  pdfauthor={\alumno}
}

% --- paleta ---
\definecolor{acento}{HTML}{1F4E79}
\definecolor{fondoIdea}{HTML}{EAF2FB}
\definecolor{fondoOjo}{HTML}{FDF0E6}
\definecolor{fondoPend}{HTML}{F2F0F7}
\definecolor{comentario}{HTML}{5A7D5A}
\definecolor{cadena}{HTML}{A03030}
\definecolor{fondoCodigo}{HTML}{F7F7F7}
\definecolor{fondoResultado}{HTML}{EAF5EC}
\definecolor{marcador}{HTML}{FFF3A3}

% --- entornos de teorema; numeran por sección, o sea por clase ---
\theoremstyle{definition}
\newtheorem{definicion}{Definición}[section]
\newtheorem{ejemplo}[definicion]{Ejemplo}
\newtheorem{ejercicio}[definicion]{Ejercicio}

\theoremstyle{plain}
\newtheorem{teorema}[definicion]{Teorema}
\newtheorem{lema}[definicion]{Lema}
\newtheorem{corolario}[definicion]{Corolario}
\newtheorem{proposicion}[definicion]{Proposición}

\theoremstyle{remark}
\newtheorem*{observacion}{Observación}
\newtheorem*{quick_sum}{Resumen rápido}

% --- cajas de color para lo que se relee antes del examen ---
\newtcolorbox{idea}[1][]{
  colback=fondoIdea, colframe=acento, boxrule=0.6pt,
  arc=2pt, left=6pt, right=6pt, top=5pt, bottom=5pt,
  title={Idea clave}, fonttitle=\bfseries\small, #1
}
\newtcolorbox{ojo}[1][]{
  colback=fondoOjo, colframe=orange!70!black, boxrule=0.6pt,
  arc=2pt, left=6pt, right=6pt, top=5pt, bottom=5pt,
  title={Ojito...}, fonttitle=\bfseries\small, #1
}
\newtcolorbox{pendiente}[1][]{
  colback=fondoPend, colframe=violet!60!black, boxrule=0.6pt,
  arc=2pt, left=6pt, right=6pt, top=5pt, bottom=5pt,
  title={Pendiente por resolver}, fonttitle=\bfseries\small, #1
}
\newtcolorbox{idear}[1][]{
  colback=fondoPend, colframe=pink!60!black, boxrule=0.6pt,
  arc=2pt, left=6pt, right=6pt, top=5pt, bottom=5pt,
  title={Idea Random}, fonttitle=\bfseries\small, #1
}

% --- el resultado al que llego un ejercicio, para que se vea de
%     lejos al hojear la hoja antes del examen ---
\newtcolorbox{resultado}[1][]{
  colback=fondoResultado, colframe=green!45!black, boxrule=0.6pt,
  arc=2pt, left=6pt, right=6pt, top=5pt, bottom=5pt,
  title={Resultado}, fonttitle=\bfseries\small, #1
}

% --- código ---
\lstdefinestyle{codigo}{
  backgroundcolor=\color{fondoCodigo},
  basicstyle=\ttfamily\footnotesize,
  keywordstyle=\color{acento}\bfseries,
  commentstyle=\color{comentario}\itshape,
  stringstyle=\color{cadena},
  numbers=left, numberstyle=\tiny\color{gray}, numbersep=8pt,
  frame=single, rulecolor=\color{gray!40},
  breaklines=true, showstringspaces=false,
  tabsize=4, captionpos=b,
  extendedchars=true, inputencoding=utf8,
  literate={á}{{\'a}}1 {é}{{\'e}}1 {í}{{\'i}}1
           {ó}{{\'o}}1 {ú}{{\'u}}1 {ñ}{{\~n}}1
           {Ü}{{\"U}}1 {ü}{{\"u}}1 {¿}{{?`}}1
           {¡}{{!`}}1
}
\lstset{style=codigo}
\renewcommand{\lstlistingname}{Código}
\renewcommand{\lstlistlistingname}{Índice de código}
% --- pseudocódigo; algorithmicx arma el cuerpo y algorithm el
%     flotante que lo numera y lo pone en el índice. Las palabras
%     clave vienen en inglés de fábrica y babel no las toca, así que
%     se renombran una por una ---
\floatname{algorithm}{Algoritmo}
\renewcommand{\listalgorithmname}{Índice de algoritmos}

\algrenewcommand\algorithmicrequire{\textbf{Entrada:}}
\algrenewcommand\algorithmicensure{\textbf{Salida:}}
\algrenewcommand\algorithmicend{\textbf{fin}}
\algrenewcommand\algorithmicif{\textbf{si}}
\algrenewcommand\algorithmicthen{\textbf{entonces}}
\algrenewcommand\algorithmicelse{\textbf{si no}}
\algrenewcommand\algorithmicfor{\textbf{para}}
\algrenewcommand\algorithmicforall{\textbf{para todo}}
\algrenewcommand\algorithmicdo{\textbf{hacer}}
\algrenewcommand\algorithmicwhile{\textbf{mientras}}
\algrenewcommand\algorithmicrepeat{\textbf{repetir}}
\algrenewcommand\algorithmicuntil{\textbf{hasta que}}
\algrenewcommand\algorithmicloop{\textbf{ciclo}}
\algrenewcommand\algorithmicprocedure{\textbf{procedimiento}}
\algrenewcommand\algorithmicfunction{\textbf{función}}
\algrenewcommand\algorithmicreturn{\textbf{devolver}}
\algrenewcomment[1]{\hfill{\color{comentario}\itshape // #1}}

% Las dos que algorithmicx no trae. El \ del final es obligatorio: sin
% el TeX se come el espacio y sale "hastan-1" en vez de "hasta n-1"
\algnewcommand\Hasta{\textbf{hasta}\ }
\algnewcommand\Intercambiar{\textbf{intercambiar}\ }

% --- abreviaturas que se usan a cada rato ---
\newcommand{\N}{\mathbb{N}}
\newcommand{\Z}{\mathbb{Z}}
\newcommand{\Q}{\mathbb{Q}}
\newcommand{\R}{\mathbb{R}}
\newcommand{\C}{\mathbb{C}}
\newcommand{\abs}[1]{\left\lvert #1 \right\rvert}
\newcommand{\norm}[1]{\left\lVert #1 \right\rVert}
\newcommand{\set}[1]{\left\{ #1 \right\}}
\newcommand{\Ord}[1]{\mathcal{O}\!\left(#1\right)}
% Mayuscula la notacion grande, minuscula la chica
\newcommand{\Omg}[1]{\Omega\!\left(#1\right)}
\newcommand{\Tht}[1]{\Theta\!\left(#1\right)}
\newcommand{\ord}[1]{o\!\left(#1\right)}
\newcommand{\omg}[1]{\omega\!\left(#1\right)}
% Los tres limites sobre n, que salen en cada renglon del criterio
% del cociente para la notacion asintotica
\newcommand{\LM}{\lim_{n \to \infty}}
\newcommand{\LS}{\limsup_{n \to \infty}}
\newcommand{\LI}{\liminf_{n \to \infty}}

% --- encabezado de cada sesión: título a la izquierda, fecha a la
%     derecha. El argumento corto es el que sale en el índice ---
\newcommand{\clase}[2]{%
  \section[#1]{#1 \hfill \normalfont\small\textit{#2}}%
}

% --- portadilla y tabla de contenidos. Se llama una vez, justo
%     despues de \begin{document}. El renglon del profesor va dentro
%     del center, asi que \profesor puede traer varios nombres
%     separados por \\ y salen apilados y centrados sin tocar nada ---
\newcommand{\portadilla}{%
  \begin{center}
    {\LARGE\bfseries\materia\par}
    \vspace{0.4em}
    {\large \profesor \par}
    \vspace{0.2em}
    {\small \periodo\quad-\quad\alumno\par}
    \vspace{0.3em}
    \rule{\linewidth}{0.5pt}
  \end{center}
  \vspace{0.5em}
  \tableofcontents
  \vspace{1em}
  \rule{\linewidth}{0.5pt}%
}

% =====================================================================
%  DOCUMENTOS DE SOLUCIONES
%  Lo que usa una tanda de ejercicios resueltos. Una libreta no toca
%  nada de aqui, pero vive en el preambulo compartido para no tener
%  dos archivos que arreglar cuando cambie el estilo.
% =====================================================================

% --- encabezado compacto: titulo, subtitulo y regla. Sustituye a
%     \portadilla cuando el documento es una tanda de ejercicios: sin
%     indice, porque se hojea de corrido, no se navega ---
\newcommand{\encabezadoSoluciones}[2]{%
  % El paquete algorithm se carga con la opcion [section] para que en
  % una libreta el pseudocodigo se numere por clase. Aqui no hay
  % secciones, asi que esa numeracion sale como "Algoritmo 0.1". Se
  % corrige al vuelo, y solo en los documentos que llaman a este
  % encabezado: la libreta no se entera.
  \renewcommand{\thealgorithm}{\arabic{algorithm}}%
  \begin{center}
    {\LARGE\bfseries #1\par}
    \vspace{0.3em}
    {\small\itshape #2\par}
    \vspace{0.5em}
    \rule{\linewidth}{0.5pt}
  \end{center}
  \vspace{0.3em}%
}

% --- hoja de un solo ejercicio: un renglon chico arriba con el
%     titulo y su numero, y ya. Sin titulo grande y sin pie de pagina,
%     porque en una hoja de dos paginas cada centimetro que se va en
%     adornos es un paso del desarrollo que no cabe ---
\newcommand{\encabezadoEjercicio}[3]{%
  % Misma correccion que en \encabezadoSoluciones: sin secciones, el
  % pseudocodigo saldria numerado "0.1"
  \renewcommand{\thealgorithm}{\arabic{algorithm}}%
  \fancyhf{}%
  \fancyhead[L]{\small\textsc{#1} (ejercicio #2 de #3)}%
  % El nombre va a la derecha porque las hojas de respuestas se
  % entregan y el profesor pide que cada una lo lleve
  \fancyhead[R]{\small\alumno}%
  \renewcommand{\headrulewidth}{0.4pt}%
  \pagestyle{fancy}%
  \thispagestyle{fancy}%
}

% --- bloque tematico, numerado con romanos: I, II, III. Agrupa los
%     ejercicios que comparten tema. No dibuja regla: la del
%     \problema que sigue ya separa el titulo del primer ejercicio,
%     y dos rayas seguidas dejan un hueco feo ---
\newcounter{parte}
\renewcommand{\theparte}{\Roman{parte}}
\newcommand{\parte}[1]{%
  \bigskip
  \refstepcounter{parte}%
  \noindent{\large\bfseries \theparte. #1\par}%
}

% --- un ejercicio resuelto. Se numera solo y corrido a lo largo del
%     documento, sin reiniciar en cada parte, para poder decir "el 14"
%     sin aclarar de que bloque ---
\newcounter{problema}
\newcommand{\problema}[1]{%
  \bigskip
  \noindent\rule{\linewidth}{0.4pt}\par
  \vspace{0.3em}
  \refstepcounter{problema}%
  \noindent{\bfseries Ejercicio \theproblema\ --- #1\par}
  \vspace{0.3em}%
}

% --- rotulo del paso que sigue: Caso base, Hipotesis inductiva, Paso
%     inductivo. El texto sigue en el mismo renglon ---
\newcommand{\paso}[1]{\par\vspace{0.3em}\noindent\textbf{#1}\ }

% --- la respuesta de un ejercicio, centrada y en un cuadro. El
%     argumento va en modo matematico, sin los delimitadores ---
\newcommand{\respuesta}[1]{%
  \[
    \boxed{\;#1\;}
  \]
}

% --- resaltado inline, como pasarle el plumon a media frase. No
%     admite matematicas adentro: soul rompe con $...$, asi que una
%     formula que hay que destacar va en \respuesta ---
\sethlcolor{marcador}
\newcommand{\resaltar}[1]{\hl{#1}}

% --- justificacion al final de un renglon de \align, en chiquito y
%     entre parentesis, para decir de donde salio el paso ---
\newcommand{\razon}[1]{\quad\text{\small\itshape (#1)}}

%############################ END OF PREAMBULO.TEX #############################
%###############################################################################

%  \begin{document}
%  \portadilla
%
%  Con dos o mas profesores se declara UN solo \profesor y los nombres
%  se separan con \\, uno por renglon:
%
%  \newcommand{\profesor}{Dra. Fulana de Tal \\
%                         Dr. Mengano de Cual}
% =====================================================================


% ==================== IDIOMA Y CODIFICACIÓN ====================
\usepackage[utf8]{inputenc}
\usepackage[T1]{fontenc}
\usepackage[spanish,es-tabla,es-noshorthands]{babel}
\usepackage{lmodern}
\usepackage{microtype}

% ==================== PÁGINA ====================
\usepackage[letterpaper,margin=2.3cm,headheight=14pt]{geometry}
\usepackage{parskip}
\usepackage{fancyhdr}
\usepackage{enumitem}

% ==================== MATEMÁTICAS ====================
\usepackage{amsmath,amssymb,amsthm,mathtools}

% ==================== FIGURAS, TABLAS, CÓDIGO ====================
\usepackage{graphicx}
\usepackage{booktabs}
\usepackage{array}
% float da el especificador [H], que clava una tabla donde se
% escribio. Una traza que LaTeX manda dos paginas adelante deja
% de leerse junto al paso que la produjo
\usepackage{float}
\usepackage{xcolor}
% tikz ya viene arrastrado por tcolorbox[most], pero se carga
% explicito para no depender de ese detalle; la libreria babel
% protege a tikz de los caracteres activos del español
\usepackage{tikz}
\usetikzlibrary{babel}
% pgfplots dibuja graficas de datos (ejes, escalas log) sobre tikz
\usepackage{pgfplots}
\pgfplotsset{compat=1.18}
\usepackage{listings}
\usepackage[ruled,section]{algorithm}
\usepackage{algpseudocode}
\usepackage[most]{tcolorbox}
% soul da \hl, el unico resaltado que parte renglon; \colorbox no
\usepackage{soul}
\usepackage{hyperref}

% =====================================================================
%  ESTILO
% =====================================================================

% --- encabezado y pie de cada página ---
\pagestyle{fancy}
\fancyhf{}
\fancyhead[L]{\small\textsc{\materia}}
\fancyhead[R]{\small\periodo}
\fancyfoot[C]{\small\thepage}
\renewcommand{\headrulewidth}{0.4pt}

% --- enlaces ---
\hypersetup{
  colorlinks=true,
  linkcolor=black,
  urlcolor=blue,
  citecolor=black,
  pdftitle={Apuntes de \materia},
  pdfauthor={\alumno}
}

% --- paleta ---
\definecolor{acento}{HTML}{1F4E79}
\definecolor{fondoIdea}{HTML}{EAF2FB}
\definecolor{fondoOjo}{HTML}{FDF0E6}
\definecolor{fondoPend}{HTML}{F2F0F7}
\definecolor{comentario}{HTML}{5A7D5A}
\definecolor{cadena}{HTML}{A03030}
\definecolor{fondoCodigo}{HTML}{F7F7F7}
\definecolor{fondoResultado}{HTML}{EAF5EC}
\definecolor{marcador}{HTML}{FFF3A3}

% --- entornos de teorema; numeran por sección, o sea por clase ---
\theoremstyle{definition}
\newtheorem{definicion}{Definición}[section]
\newtheorem{ejemplo}[definicion]{Ejemplo}
\newtheorem{ejercicio}[definicion]{Ejercicio}

\theoremstyle{plain}
\newtheorem{teorema}[definicion]{Teorema}
\newtheorem{lema}[definicion]{Lema}
\newtheorem{corolario}[definicion]{Corolario}
\newtheorem{proposicion}[definicion]{Proposición}

\theoremstyle{remark}
\newtheorem*{observacion}{Observación}
\newtheorem*{quick_sum}{Resumen rápido}

% --- cajas de color para lo que se relee antes del examen ---
\newtcolorbox{idea}[1][]{
  colback=fondoIdea, colframe=acento, boxrule=0.6pt,
  arc=2pt, left=6pt, right=6pt, top=5pt, bottom=5pt,
  title={Idea clave}, fonttitle=\bfseries\small, #1
}
\newtcolorbox{ojo}[1][]{
  colback=fondoOjo, colframe=orange!70!black, boxrule=0.6pt,
  arc=2pt, left=6pt, right=6pt, top=5pt, bottom=5pt,
  title={Ojito...}, fonttitle=\bfseries\small, #1
}
\newtcolorbox{pendiente}[1][]{
  colback=fondoPend, colframe=violet!60!black, boxrule=0.6pt,
  arc=2pt, left=6pt, right=6pt, top=5pt, bottom=5pt,
  title={Pendiente por resolver}, fonttitle=\bfseries\small, #1
}
\newtcolorbox{idear}[1][]{
  colback=fondoPend, colframe=pink!60!black, boxrule=0.6pt,
  arc=2pt, left=6pt, right=6pt, top=5pt, bottom=5pt,
  title={Idea Random}, fonttitle=\bfseries\small, #1
}

% --- el resultado al que llego un ejercicio, para que se vea de
%     lejos al hojear la hoja antes del examen ---
\newtcolorbox{resultado}[1][]{
  colback=fondoResultado, colframe=green!45!black, boxrule=0.6pt,
  arc=2pt, left=6pt, right=6pt, top=5pt, bottom=5pt,
  title={Resultado}, fonttitle=\bfseries\small, #1
}

% --- código ---
\lstdefinestyle{codigo}{
  backgroundcolor=\color{fondoCodigo},
  basicstyle=\ttfamily\footnotesize,
  keywordstyle=\color{acento}\bfseries,
  commentstyle=\color{comentario}\itshape,
  stringstyle=\color{cadena},
  numbers=left, numberstyle=\tiny\color{gray}, numbersep=8pt,
  frame=single, rulecolor=\color{gray!40},
  breaklines=true, showstringspaces=false,
  tabsize=4, captionpos=b,
  extendedchars=true, inputencoding=utf8,
  literate={á}{{\'a}}1 {é}{{\'e}}1 {í}{{\'i}}1
           {ó}{{\'o}}1 {ú}{{\'u}}1 {ñ}{{\~n}}1
           {Ü}{{\"U}}1 {ü}{{\"u}}1 {¿}{{?`}}1
           {¡}{{!`}}1
}
\lstset{style=codigo}
\renewcommand{\lstlistingname}{Código}
\renewcommand{\lstlistlistingname}{Índice de código}
% --- pseudocódigo; algorithmicx arma el cuerpo y algorithm el
%     flotante que lo numera y lo pone en el índice. Las palabras
%     clave vienen en inglés de fábrica y babel no las toca, así que
%     se renombran una por una ---
\floatname{algorithm}{Algoritmo}
\renewcommand{\listalgorithmname}{Índice de algoritmos}

\algrenewcommand\algorithmicrequire{\textbf{Entrada:}}
\algrenewcommand\algorithmicensure{\textbf{Salida:}}
\algrenewcommand\algorithmicend{\textbf{fin}}
\algrenewcommand\algorithmicif{\textbf{si}}
\algrenewcommand\algorithmicthen{\textbf{entonces}}
\algrenewcommand\algorithmicelse{\textbf{si no}}
\algrenewcommand\algorithmicfor{\textbf{para}}
\algrenewcommand\algorithmicforall{\textbf{para todo}}
\algrenewcommand\algorithmicdo{\textbf{hacer}}
\algrenewcommand\algorithmicwhile{\textbf{mientras}}
\algrenewcommand\algorithmicrepeat{\textbf{repetir}}
\algrenewcommand\algorithmicuntil{\textbf{hasta que}}
\algrenewcommand\algorithmicloop{\textbf{ciclo}}
\algrenewcommand\algorithmicprocedure{\textbf{procedimiento}}
\algrenewcommand\algorithmicfunction{\textbf{función}}
\algrenewcommand\algorithmicreturn{\textbf{devolver}}
\algrenewcomment[1]{\hfill{\color{comentario}\itshape // #1}}

% Las dos que algorithmicx no trae. El \ del final es obligatorio: sin
% el TeX se come el espacio y sale "hastan-1" en vez de "hasta n-1"
\algnewcommand\Hasta{\textbf{hasta}\ }
\algnewcommand\Intercambiar{\textbf{intercambiar}\ }

% --- abreviaturas que se usan a cada rato ---
\newcommand{\N}{\mathbb{N}}
\newcommand{\Z}{\mathbb{Z}}
\newcommand{\Q}{\mathbb{Q}}
\newcommand{\R}{\mathbb{R}}
\newcommand{\C}{\mathbb{C}}
\newcommand{\abs}[1]{\left\lvert #1 \right\rvert}
\newcommand{\norm}[1]{\left\lVert #1 \right\rVert}
\newcommand{\set}[1]{\left\{ #1 \right\}}
\newcommand{\Ord}[1]{\mathcal{O}\!\left(#1\right)}
% Mayuscula la notacion grande, minuscula la chica
\newcommand{\Omg}[1]{\Omega\!\left(#1\right)}
\newcommand{\Tht}[1]{\Theta\!\left(#1\right)}
\newcommand{\ord}[1]{o\!\left(#1\right)}
\newcommand{\omg}[1]{\omega\!\left(#1\right)}
% Los tres limites sobre n, que salen en cada renglon del criterio
% del cociente para la notacion asintotica
\newcommand{\LM}{\lim_{n \to \infty}}
\newcommand{\LS}{\limsup_{n \to \infty}}
\newcommand{\LI}{\liminf_{n \to \infty}}

% --- encabezado de cada sesión: título a la izquierda, fecha a la
%     derecha. El argumento corto es el que sale en el índice ---
\newcommand{\clase}[2]{%
  \section[#1]{#1 \hfill \normalfont\small\textit{#2}}%
}

% --- portadilla y tabla de contenidos. Se llama una vez, justo
%     despues de \begin{document}. El renglon del profesor va dentro
%     del center, asi que \profesor puede traer varios nombres
%     separados por \\ y salen apilados y centrados sin tocar nada ---
\newcommand{\portadilla}{%
  \begin{center}
    {\LARGE\bfseries\materia\par}
    \vspace{0.4em}
    {\large \profesor \par}
    \vspace{0.2em}
    {\small \periodo\quad-\quad\alumno\par}
    \vspace{0.3em}
    \rule{\linewidth}{0.5pt}
  \end{center}
  \vspace{0.5em}
  \tableofcontents
  \vspace{1em}
  \rule{\linewidth}{0.5pt}%
}

% =====================================================================
%  DOCUMENTOS DE SOLUCIONES
%  Lo que usa una tanda de ejercicios resueltos. Una libreta no toca
%  nada de aqui, pero vive en el preambulo compartido para no tener
%  dos archivos que arreglar cuando cambie el estilo.
% =====================================================================

% --- encabezado compacto: titulo, subtitulo y regla. Sustituye a
%     \portadilla cuando el documento es una tanda de ejercicios: sin
%     indice, porque se hojea de corrido, no se navega ---
\newcommand{\encabezadoSoluciones}[2]{%
  % El paquete algorithm se carga con la opcion [section] para que en
  % una libreta el pseudocodigo se numere por clase. Aqui no hay
  % secciones, asi que esa numeracion sale como "Algoritmo 0.1". Se
  % corrige al vuelo, y solo en los documentos que llaman a este
  % encabezado: la libreta no se entera.
  \renewcommand{\thealgorithm}{\arabic{algorithm}}%
  \begin{center}
    {\LARGE\bfseries #1\par}
    \vspace{0.3em}
    {\small\itshape #2\par}
    \vspace{0.5em}
    \rule{\linewidth}{0.5pt}
  \end{center}
  \vspace{0.3em}%
}

% --- hoja de un solo ejercicio: un renglon chico arriba con el
%     titulo y su numero, y ya. Sin titulo grande y sin pie de pagina,
%     porque en una hoja de dos paginas cada centimetro que se va en
%     adornos es un paso del desarrollo que no cabe ---
\newcommand{\encabezadoEjercicio}[3]{%
  % Misma correccion que en \encabezadoSoluciones: sin secciones, el
  % pseudocodigo saldria numerado "0.1"
  \renewcommand{\thealgorithm}{\arabic{algorithm}}%
  \fancyhf{}%
  \fancyhead[L]{\small\textsc{#1} (ejercicio #2 de #3)}%
  % El nombre va a la derecha porque las hojas de respuestas se
  % entregan y el profesor pide que cada una lo lleve
  \fancyhead[R]{\small\alumno}%
  \renewcommand{\headrulewidth}{0.4pt}%
  \pagestyle{fancy}%
  \thispagestyle{fancy}%
}

% --- bloque tematico, numerado con romanos: I, II, III. Agrupa los
%     ejercicios que comparten tema. No dibuja regla: la del
%     \problema que sigue ya separa el titulo del primer ejercicio,
%     y dos rayas seguidas dejan un hueco feo ---
\newcounter{parte}
\renewcommand{\theparte}{\Roman{parte}}
\newcommand{\parte}[1]{%
  \bigskip
  \refstepcounter{parte}%
  \noindent{\large\bfseries \theparte. #1\par}%
}

% --- un ejercicio resuelto. Se numera solo y corrido a lo largo del
%     documento, sin reiniciar en cada parte, para poder decir "el 14"
%     sin aclarar de que bloque ---
\newcounter{problema}
\newcommand{\problema}[1]{%
  \bigskip
  \noindent\rule{\linewidth}{0.4pt}\par
  \vspace{0.3em}
  \refstepcounter{problema}%
  \noindent{\bfseries Ejercicio \theproblema\ --- #1\par}
  \vspace{0.3em}%
}

% --- rotulo del paso que sigue: Caso base, Hipotesis inductiva, Paso
%     inductivo. El texto sigue en el mismo renglon ---
\newcommand{\paso}[1]{\par\vspace{0.3em}\noindent\textbf{#1}\ }

% --- la respuesta de un ejercicio, centrada y en un cuadro. El
%     argumento va en modo matematico, sin los delimitadores ---
\newcommand{\respuesta}[1]{%
  \[
    \boxed{\;#1\;}
  \]
}

% --- resaltado inline, como pasarle el plumon a media frase. No
%     admite matematicas adentro: soul rompe con $...$, asi que una
%     formula que hay que destacar va en \respuesta ---
\sethlcolor{marcador}
\newcommand{\resaltar}[1]{\hl{#1}}

% --- justificacion al final de un renglon de \align, en chiquito y
%     entre parentesis, para decir de donde salio el paso ---
\newcommand{\razon}[1]{\quad\text{\small\itshape (#1)}}

%############################ END OF PREAMBULO.TEX #############################
%###############################################################################

%  \begin{document}
%  \portadilla
%
%  Con dos o mas profesores se declara UN solo \profesor y los nombres
%  se separan con \\, uno por renglon:
%
%  \newcommand{\profesor}{Dra. Fulana de Tal \\
%                         Dr. Mengano de Cual}
% =====================================================================


% ==================== IDIOMA Y CODIFICACIÓN ====================
\usepackage[utf8]{inputenc}
\usepackage[T1]{fontenc}
\usepackage[spanish,es-tabla,es-noshorthands]{babel}
\usepackage{lmodern}
\usepackage{microtype}

% ==================== PÁGINA ====================
\usepackage[letterpaper,margin=2.3cm,headheight=14pt]{geometry}
\usepackage{parskip}
\usepackage{fancyhdr}
\usepackage{enumitem}

% ==================== MATEMÁTICAS ====================
\usepackage{amsmath,amssymb,amsthm,mathtools}

% ==================== FIGURAS, TABLAS, CÓDIGO ====================
\usepackage{graphicx}
\usepackage{booktabs}
\usepackage{array}
% float da el especificador [H], que clava una tabla donde se
% escribio. Una traza que LaTeX manda dos paginas adelante deja
% de leerse junto al paso que la produjo
\usepackage{float}
\usepackage{xcolor}
% tikz ya viene arrastrado por tcolorbox[most], pero se carga
% explicito para no depender de ese detalle; la libreria babel
% protege a tikz de los caracteres activos del español
\usepackage{tikz}
\usetikzlibrary{babel}
% pgfplots dibuja graficas de datos (ejes, escalas log) sobre tikz
\usepackage{pgfplots}
\pgfplotsset{compat=1.18}
\usepackage{listings}
\usepackage[ruled,section]{algorithm}
\usepackage{algpseudocode}
\usepackage[most]{tcolorbox}
% soul da \hl, el unico resaltado que parte renglon; \colorbox no
\usepackage{soul}
\usepackage{hyperref}

% =====================================================================
%  ESTILO
% =====================================================================

% --- encabezado y pie de cada página ---
\pagestyle{fancy}
\fancyhf{}
\fancyhead[L]{\small\textsc{\materia}}
\fancyhead[R]{\small\periodo}
\fancyfoot[C]{\small\thepage}
\renewcommand{\headrulewidth}{0.4pt}

% --- enlaces ---
\hypersetup{
  colorlinks=true,
  linkcolor=black,
  urlcolor=blue,
  citecolor=black,
  pdftitle={Apuntes de \materia},
  pdfauthor={\alumno}
}

% --- paleta ---
\definecolor{acento}{HTML}{1F4E79}
\definecolor{fondoIdea}{HTML}{EAF2FB}
\definecolor{fondoOjo}{HTML}{FDF0E6}
\definecolor{fondoPend}{HTML}{F2F0F7}
\definecolor{comentario}{HTML}{5A7D5A}
\definecolor{cadena}{HTML}{A03030}
\definecolor{fondoCodigo}{HTML}{F7F7F7}
\definecolor{fondoResultado}{HTML}{EAF5EC}
\definecolor{marcador}{HTML}{FFF3A3}

% --- entornos de teorema; numeran por sección, o sea por clase ---
\theoremstyle{definition}
\newtheorem{definicion}{Definición}[section]
\newtheorem{ejemplo}[definicion]{Ejemplo}
\newtheorem{ejercicio}[definicion]{Ejercicio}

\theoremstyle{plain}
\newtheorem{teorema}[definicion]{Teorema}
\newtheorem{lema}[definicion]{Lema}
\newtheorem{corolario}[definicion]{Corolario}
\newtheorem{proposicion}[definicion]{Proposición}

\theoremstyle{remark}
\newtheorem*{observacion}{Observación}
\newtheorem*{quick_sum}{Resumen rápido}

% --- cajas de color para lo que se relee antes del examen ---
\newtcolorbox{idea}[1][]{
  colback=fondoIdea, colframe=acento, boxrule=0.6pt,
  arc=2pt, left=6pt, right=6pt, top=5pt, bottom=5pt,
  title={Idea clave}, fonttitle=\bfseries\small, #1
}
\newtcolorbox{ojo}[1][]{
  colback=fondoOjo, colframe=orange!70!black, boxrule=0.6pt,
  arc=2pt, left=6pt, right=6pt, top=5pt, bottom=5pt,
  title={Ojito...}, fonttitle=\bfseries\small, #1
}
\newtcolorbox{pendiente}[1][]{
  colback=fondoPend, colframe=violet!60!black, boxrule=0.6pt,
  arc=2pt, left=6pt, right=6pt, top=5pt, bottom=5pt,
  title={Pendiente por resolver}, fonttitle=\bfseries\small, #1
}
\newtcolorbox{idear}[1][]{
  colback=fondoPend, colframe=pink!60!black, boxrule=0.6pt,
  arc=2pt, left=6pt, right=6pt, top=5pt, bottom=5pt,
  title={Idea Random}, fonttitle=\bfseries\small, #1
}

% --- el resultado al que llego un ejercicio, para que se vea de
%     lejos al hojear la hoja antes del examen ---
\newtcolorbox{resultado}[1][]{
  colback=fondoResultado, colframe=green!45!black, boxrule=0.6pt,
  arc=2pt, left=6pt, right=6pt, top=5pt, bottom=5pt,
  title={Resultado}, fonttitle=\bfseries\small, #1
}

% --- código ---
\lstdefinestyle{codigo}{
  backgroundcolor=\color{fondoCodigo},
  basicstyle=\ttfamily\footnotesize,
  keywordstyle=\color{acento}\bfseries,
  commentstyle=\color{comentario}\itshape,
  stringstyle=\color{cadena},
  numbers=left, numberstyle=\tiny\color{gray}, numbersep=8pt,
  frame=single, rulecolor=\color{gray!40},
  breaklines=true, showstringspaces=false,
  tabsize=4, captionpos=b,
  extendedchars=true, inputencoding=utf8,
  literate={á}{{\'a}}1 {é}{{\'e}}1 {í}{{\'i}}1
           {ó}{{\'o}}1 {ú}{{\'u}}1 {ñ}{{\~n}}1
           {Ü}{{\"U}}1 {ü}{{\"u}}1 {¿}{{?`}}1
           {¡}{{!`}}1
}
\lstset{style=codigo}
\renewcommand{\lstlistingname}{Código}
\renewcommand{\lstlistlistingname}{Índice de código}
% --- pseudocódigo; algorithmicx arma el cuerpo y algorithm el
%     flotante que lo numera y lo pone en el índice. Las palabras
%     clave vienen en inglés de fábrica y babel no las toca, así que
%     se renombran una por una ---
\floatname{algorithm}{Algoritmo}
\renewcommand{\listalgorithmname}{Índice de algoritmos}

\algrenewcommand\algorithmicrequire{\textbf{Entrada:}}
\algrenewcommand\algorithmicensure{\textbf{Salida:}}
\algrenewcommand\algorithmicend{\textbf{fin}}
\algrenewcommand\algorithmicif{\textbf{si}}
\algrenewcommand\algorithmicthen{\textbf{entonces}}
\algrenewcommand\algorithmicelse{\textbf{si no}}
\algrenewcommand\algorithmicfor{\textbf{para}}
\algrenewcommand\algorithmicforall{\textbf{para todo}}
\algrenewcommand\algorithmicdo{\textbf{hacer}}
\algrenewcommand\algorithmicwhile{\textbf{mientras}}
\algrenewcommand\algorithmicrepeat{\textbf{repetir}}
\algrenewcommand\algorithmicuntil{\textbf{hasta que}}
\algrenewcommand\algorithmicloop{\textbf{ciclo}}
\algrenewcommand\algorithmicprocedure{\textbf{procedimiento}}
\algrenewcommand\algorithmicfunction{\textbf{función}}
\algrenewcommand\algorithmicreturn{\textbf{devolver}}
\algrenewcomment[1]{\hfill{\color{comentario}\itshape // #1}}

% Las dos que algorithmicx no trae. El \ del final es obligatorio: sin
% el TeX se come el espacio y sale "hastan-1" en vez de "hasta n-1"
\algnewcommand\Hasta{\textbf{hasta}\ }
\algnewcommand\Intercambiar{\textbf{intercambiar}\ }

% --- abreviaturas que se usan a cada rato ---
\newcommand{\N}{\mathbb{N}}
\newcommand{\Z}{\mathbb{Z}}
\newcommand{\Q}{\mathbb{Q}}
\newcommand{\R}{\mathbb{R}}
\newcommand{\C}{\mathbb{C}}
\newcommand{\abs}[1]{\left\lvert #1 \right\rvert}
\newcommand{\norm}[1]{\left\lVert #1 \right\rVert}
\newcommand{\set}[1]{\left\{ #1 \right\}}
\newcommand{\Ord}[1]{\mathcal{O}\!\left(#1\right)}
% Mayuscula la notacion grande, minuscula la chica
\newcommand{\Omg}[1]{\Omega\!\left(#1\right)}
\newcommand{\Tht}[1]{\Theta\!\left(#1\right)}
\newcommand{\ord}[1]{o\!\left(#1\right)}
\newcommand{\omg}[1]{\omega\!\left(#1\right)}
% Los tres limites sobre n, que salen en cada renglon del criterio
% del cociente para la notacion asintotica
\newcommand{\LM}{\lim_{n \to \infty}}
\newcommand{\LS}{\limsup_{n \to \infty}}
\newcommand{\LI}{\liminf_{n \to \infty}}

% --- encabezado de cada sesión: título a la izquierda, fecha a la
%     derecha. El argumento corto es el que sale en el índice ---
\newcommand{\clase}[2]{%
  \section[#1]{#1 \hfill \normalfont\small\textit{#2}}%
}

% --- portadilla y tabla de contenidos. Se llama una vez, justo
%     despues de \begin{document}. El renglon del profesor va dentro
%     del center, asi que \profesor puede traer varios nombres
%     separados por \\ y salen apilados y centrados sin tocar nada ---
\newcommand{\portadilla}{%
  \begin{center}
    {\LARGE\bfseries\materia\par}
    \vspace{0.4em}
    {\large \profesor \par}
    \vspace{0.2em}
    {\small \periodo\quad-\quad\alumno\par}
    \vspace{0.3em}
    \rule{\linewidth}{0.5pt}
  \end{center}
  \vspace{0.5em}
  \tableofcontents
  \vspace{1em}
  \rule{\linewidth}{0.5pt}%
}

% =====================================================================
%  DOCUMENTOS DE SOLUCIONES
%  Lo que usa una tanda de ejercicios resueltos. Una libreta no toca
%  nada de aqui, pero vive en el preambulo compartido para no tener
%  dos archivos que arreglar cuando cambie el estilo.
% =====================================================================

% --- encabezado compacto: titulo, subtitulo y regla. Sustituye a
%     \portadilla cuando el documento es una tanda de ejercicios: sin
%     indice, porque se hojea de corrido, no se navega ---
\newcommand{\encabezadoSoluciones}[2]{%
  % El paquete algorithm se carga con la opcion [section] para que en
  % una libreta el pseudocodigo se numere por clase. Aqui no hay
  % secciones, asi que esa numeracion sale como "Algoritmo 0.1". Se
  % corrige al vuelo, y solo en los documentos que llaman a este
  % encabezado: la libreta no se entera.
  \renewcommand{\thealgorithm}{\arabic{algorithm}}%
  \begin{center}
    {\LARGE\bfseries #1\par}
    \vspace{0.3em}
    {\small\itshape #2\par}
    \vspace{0.5em}
    \rule{\linewidth}{0.5pt}
  \end{center}
  \vspace{0.3em}%
}

% --- hoja de un solo ejercicio: un renglon chico arriba con el
%     titulo y su numero, y ya. Sin titulo grande y sin pie de pagina,
%     porque en una hoja de dos paginas cada centimetro que se va en
%     adornos es un paso del desarrollo que no cabe ---
\newcommand{\encabezadoEjercicio}[3]{%
  % Misma correccion que en \encabezadoSoluciones: sin secciones, el
  % pseudocodigo saldria numerado "0.1"
  \renewcommand{\thealgorithm}{\arabic{algorithm}}%
  \fancyhf{}%
  \fancyhead[L]{\small\textsc{#1} (ejercicio #2 de #3)}%
  % El nombre va a la derecha porque las hojas de respuestas se
  % entregan y el profesor pide que cada una lo lleve
  \fancyhead[R]{\small\alumno}%
  \renewcommand{\headrulewidth}{0.4pt}%
  \pagestyle{fancy}%
  \thispagestyle{fancy}%
}

% --- bloque tematico, numerado con romanos: I, II, III. Agrupa los
%     ejercicios que comparten tema. No dibuja regla: la del
%     \problema que sigue ya separa el titulo del primer ejercicio,
%     y dos rayas seguidas dejan un hueco feo ---
\newcounter{parte}
\renewcommand{\theparte}{\Roman{parte}}
\newcommand{\parte}[1]{%
  \bigskip
  \refstepcounter{parte}%
  \noindent{\large\bfseries \theparte. #1\par}%
}

% --- un ejercicio resuelto. Se numera solo y corrido a lo largo del
%     documento, sin reiniciar en cada parte, para poder decir "el 14"
%     sin aclarar de que bloque ---
\newcounter{problema}
\newcommand{\problema}[1]{%
  \bigskip
  \noindent\rule{\linewidth}{0.4pt}\par
  \vspace{0.3em}
  \refstepcounter{problema}%
  \noindent{\bfseries Ejercicio \theproblema\ --- #1\par}
  \vspace{0.3em}%
}

% --- rotulo del paso que sigue: Caso base, Hipotesis inductiva, Paso
%     inductivo. El texto sigue en el mismo renglon ---
\newcommand{\paso}[1]{\par\vspace{0.3em}\noindent\textbf{#1}\ }

% --- la respuesta de un ejercicio, centrada y en un cuadro. El
%     argumento va en modo matematico, sin los delimitadores ---
\newcommand{\respuesta}[1]{%
  \[
    \boxed{\;#1\;}
  \]
}

% --- resaltado inline, como pasarle el plumon a media frase. No
%     admite matematicas adentro: soul rompe con $...$, asi que una
%     formula que hay que destacar va en \respuesta ---
\sethlcolor{marcador}
\newcommand{\resaltar}[1]{\hl{#1}}

% --- justificacion al final de un renglon de \align, en chiquito y
%     entre parentesis, para decir de donde salio el paso ---
\newcommand{\razon}[1]{\quad\text{\small\itshape (#1)}}

%############################ END OF PREAMBULO.TEX #############################
%###############################################################################


% =====================================================================
\begin{document}
% =====================================================================

\portadilla
\newpage

% =====================================================================
%  Clase 1 - Intro al curso
% =====================================================================

\clase{Introducción al Curso}{25 de agosto de 2026}

\begin{quick_sum}
  Sesión inicial: contenido y dinámica del curso, temario, criterios
  de evaluación y el calendario del proyecto. La segunda mitad fue la
  introducción propiamente dicha: qué es aprender, por qué generalizar
  no es memorizar, de qué disciplinas se alimenta el área y la
  formalización de la función $f$ que se quiere aprender contra la
  hipótesis $h$.
\end{quick_sum}

\subsection*{Info general}

\textbf{Contacto:}
\begin{itemize}[itemsep=0.1em]
    \item \profesor
    \item Correo: \texttt{fmartine@inaoep.mx}
\end{itemize}

\textbf{Temario:}
\begin{enumerate}[itemsep=0.1em]
    \item Introducción y clasificadores lineales
    \item Perceptrón
    \item SVM
    \item Clasificadores basados en distancia
    \item ANNs
    \item Evaluación, validación y overfitting
    \item Métodos bayesianos
    \item Árboles de decisión
    \item Reglas de clasificación
    \item Reglas de asociación
    \item Selección de atributos
    \item Clustering (K-means, gaussian mixture, jerárquicos)
\end{enumerate}

\begin{idea}
  No se trata de usar las herramientas como una caja negra. El
  objetivo es, como investigadores, saber cómo usarlas para generar
  nuevo conocimiento, y por eso hay que conocerlas a detalle: para
  poder innovar y proponer alternativas.
\end{idea}

% =====================================================================

\subsection*{Evaluación}

\textbf{80\% proyecto:}
\begin{itemize}[itemsep=0.1em]
    \item Equipo de dos personas
    \item Seleccionar un artículo reciente
    \item Proponer una solución alternativa
    \item Tres presentaciones
    \item Reporte final
\end{itemize}

\textbf{20\% examen oral}, a la entrega del proyecto, sobre un tema
aleatorio del temario.

% =====================================================================

\subsection*{Calendario}

\begin{itemize}[itemsep=0.4em]
    \item \textbf{25 de agosto.} Buscar un artículo publicado en 2026
          que sea de interés, en una revista JCR.
    \item \textbf{3 de septiembre.} Entregar por correo una lista que
          indique cuál artículo seleccionó cada estudiante o equipo.
    \item \textbf{29 de septiembre o 1 de octubre.} Primera
          presentación:
          \begin{itemize}[itemsep=0.1em]
              \item Descripción del problema abordado en el artículo
              \item Breve descripción del trabajo relacionado que
                    menciona el artículo
              \item Breve descripción de la solución propuesta en el
                    artículo
              \item Descripción de los experimentos que soportan dicha
                    solución, incluyendo tablas y gráficas tomadas del
                    artículo
              \item Conclusiones a las que llegan los autores
              \item Análisis de debilidades y posibles modificaciones
                    para mejorar la propuesta de los autores
          \end{itemize}
    \item \textbf{27 o 29 de octubre.} Segunda presentación: resumen
          del problema abordado.
    \item \textbf{24 o 26 de noviembre.} Presentación final: resumen
          del problema abordado y de la solución final.
    \item \textbf{30 de noviembre.} Entrega del reporte final. Se
          entrega en forma de artículo, con formato de
          \url{https://www.springer.com/}.
\end{itemize}

El reporte debe tener más o menos 12 páginas e incluir título,
resumen, introducción, trabajo relacionado, solución propuesta,
experimentos, conclusiones y trabajo futuro, y referencias.

% =====================================================================

\subsection*{Revistas JCR}

\begin{itemize}[itemsep=0.1em]
    \item Expert Systems with Applications
    \item Knowledge-Based Systems
    \item Journal of Machine Learning Research
\end{itemize}

Buscador: \url{https://www.webofscience.com/wos/woscc/basic-research}

\subsection*{Recursos}

\begin{itemize}[itemsep=0.1em]
    \item Weka 3
    \item scikit-learn
\end{itemize}

\begin{ojo}
  Hay que tratar de contactar al autor del artículo para que nos dé
  sus recursos y poder compararnos.
\end{ojo}

\newpage

% =====================================================================
%  Intro al aprendizaje
% =====================================================================

\section{Aprendizaje}

\subsection*{Que es el aprendizaje?}

Es una disciplina dentro de las ciencias de la computación que surge
como motivación de hacer que las computadoras realicen tareas
complejas como las que hacen los humanos.

Aprender significa adquirir conocimiento, entendimiento o una
habilidad de manera empírica o por estudio. La pregunta que queda
abierta es cómo determinamos el éxito o el fracaso de la tarea.

Se busca que se logre un proceso de aprendizaje para que el sistema
sea capaz de generalizar.

\begin{idea}
  Generalizar no es memorizar. Aprender permite responder a problemas
  nuevos con el conocimiento adquirido.
\end{idea}

Al inicio hubo varias tendencias. Una de ellas fue ver cómo aprenden
los animales y de ahí trasladarlo a la máquina, con la premisa de que
quizá el aprendizaje en un animal es diferente al de un ser humano,
pero no resultó cierto.

\begin{ojo}
  No hay una demostración universal de lo que es el aprendizaje. Hay
  muchas definiciones y muchas escuelas.
\end{ojo}

\subsubsection*{Tipos de aprendizaje}

\begin{itemize}[itemsep=0.1em]
    \item Supervisado
    \item No supervisado
    \item Por refuerzo
\end{itemize}

\subsubsection*{Tareas relacionadas con el aprendizaje}

\begin{itemize}[itemsep=0.1em]
    \item Reconocimiento
    \item Diagnóstico
    \item Planeación
    \item Control de robots
    \item Predicción
\end{itemize}

Aprender puede ser percibir cambios en el comportamiento gracias a la
experiencia. En nuestro contexto serían mejoras en actividades ya
conocidas, y aprender tareas nuevas.

Hacer que las máquinas aprendan es muy útil, ya que nos simplifica la
vida en muchos contextos: nos ayuda a hacer tareas complejas más
rápido y a analizar datos masivos de una manera rápida.

% =====================================================================

\subsection*{Disciplinas relacionadas}

\begin{itemize}[itemsep=0.1em]
    \item Estadística
    \item Lógica
    \item Matemáticas discretas
    \item Y muchas más
\end{itemize}

Todas las teorías que hay sobre cómo funciona el cerebro han ayudado
bastante, ya que se trata de emular lo que se cree que ocurre en el
cerebro humano cuando ocurre el aprendizaje, y trasladarlo a la
computadora.

Otra disciplina que se vincula mucho con el ML es el control, la
teoría de control. Igual los modelos psicológicos y los modelos
evolutivos, entre muchos más.

% =====================================================================

\subsection*{Qué es lo que se aprende?}

Diferentes estructuras computacionales: funciones, programas lógicos,
máquinas de estado finito, gramáticas y sistemas para la resolución de
problemas.

La idea es que tengo mi problema y tengo algunos ejemplos de él. Le
doy a la computadora los ejemplos y la computadora me da como
resultado el programa que me ayuda a generalizar cómo son esos
ejemplos, de tal manera que cuando llegue un dato nuevo pueda
clasificarlo, identificarlo y reconocerlo correctamente.

Más formalmente, se tiene una función $f$, que es la que queremos
aprender, y otra función $h$, la hipótesis. Ambas reciben un vector de
valores con $N$ componentes. Queremos encontrar la función $h$ cuya
salida sea $h(X)$, y que sea capaz de reconocer un dato nuevo.

\subsubsection*{Supervisado y no supervisado}

Hay dos tipos de aprendizaje, bueno tres, pero estos son los más
comunes.

\begin{itemize}[itemsep=0.3em]
    \item \textbf{Supervisado.} Los datos de entrenamiento están
          etiquetados. A la función $h$, que es la que voy a aprender,
          le voy a dar un dato $x$, y esa función tiene que ser capaz
          de reconocer la nueva entrada.
    \item \textbf{No supervisado.} Los datos de entrenamiento no están
          etiquetados, y el algoritmo encuentra la relación entre los
          datos.
\end{itemize}

\subsubsection*{Vector de entrada}

Al vector de entrada se le llama patrón, feature vector, instancia o
ejemplo. Sus componentes son las variables o características, y esas
variables pueden tomar valores reales, discretos o categóricos.

La salida puede ser un valor real, y entonces se trata de un
predictor, o discreto, y entonces se trata de un clasificador.

\subsubsection*{Modos de entrenamiento}

Hay varios modos de entrenamiento, dependiendo de los datos:

\begin{itemize}[itemsep=0.1em]
    \item Por batch
    \item De manera incremental
    \item En línea
\end{itemize}

% =====================================================================

\subsection*{Evaluación}

Cuando ya tengo mi función $h$, ¿cómo la evalúo? ¿Cómo sé que está
bien? Para eso se necesita un conjunto de prueba, que nos ayuda a ver
qué tan bien la función $h$ realiza su tarea, midiendo el error de $h$
sobre ese conjunto. El conjunto de prueba es una partición del dataset
original.

\newpage

% =====================================================================
%  Clase 2 - Clasificadores lineales y perceptrón
% =====================================================================

\clase{Clasificadores lineales y perceptrón}{27 de agosto de 2026}

\begin{quick_sum}
  Retomamos la idea de aprender una función que generalice y la
  aterrizamos en un ejemplo de control de acceso por rostro. De ahí
  salieron los clasificadores lineales, las funciones de pérdida y el
  perceptrón como algoritmo para ajustar los parámetros.
\end{quick_sum}

\subsection*{Recap}

ML se trata de encontrar una función que fitee unos datos de
entrenamiento para que sea capaz de generalizar. El proceso es:

\begin{enumerate}[itemsep=0.1em]
    \item Suponemos una función $f$ que modela lo que queremos
    \item Construimos hipótesis
    \item Queremos construir la función $h$ que emule a $f$
    \item Entrenamos a $h$ para que sea capaz de generalizar
    \item Testeamos
\end{enumerate}

% =====================================================================

\subsection*{Ejemplo: control de acceso por rostro}

Supongamos que queremos construir un sistema que permita pasar a
personas cuyos rostros sean conocidos para entrar al instituto.
Tenemos los datos de las caras, y una imagen es un arreglo de bits.

Tenemos que construir una función
\[
    f: \R^d \to \set{-1, 1}
\]

Es discreta porque es una tarea de clasificación. El clasificador
tendrá un conjunto de vectores de entrenamiento con etiquetas, y estos
vectores nos ayudarán a entrenar la función. Entonces queremos
construir la función solo con los datos de entrenamiento, o sea con
nuestra experiencia empírica.

Una manera sería ver un píxel de la imagen, de tal forma que al estar
presente acepta a esa persona y la rechaza en otro caso, pero esto
tiene varias limitaciones. ¿La función es útil para otras imágenes que
no estén en el dataset? ¿Es capaz de generalizar?

\begin{idea}
  Lo que se trata es de encontrar un modelo, una función que nos
  resuelva esa tarea de aprendizaje. No queremos funciones que fiteen
  perfectamente al conjunto de datos, ni súper generales, queremos
  algo en medio.
\end{idea}

\subsubsection*{Selección de modelo}

Y cómo se encuentran clasificadores que generalicen bien. La clave es
\emph{restringir} el conjunto de funciones binarias que uno se permite
considerar. Se busca una clase de funciones tal que, si una función de
la clase funciona bien en el entrenamiento, es probable que también
funcione bien en las imágenes que no se han visto.

\begin{idea}
  \textbf{El problema de selección de modelo.} La clase correcta no
  puede ser demasiado grande, porque entonces contiene demasiadas
  funciones claramente distintas y una de ellas va a memorizar el
  entrenamiento por casualidad, como la del píxel único. Pero tampoco
  demasiado chica, porque entonces se corre el riesgo de no encontrar
  ninguna función que funcione bien ni siquiera en el entrenamiento.
  Encontrar esa clase es uno de los problemas centrales del
  aprendizaje automático.
\end{idea}

Fijar la clase de funciones en los clasificadores lineales, que es lo
que sigue, es exactamente tomar esa decisión.

% =====================================================================

\subsection*{Clasificadores lineales}

Lo que hacemos es definir una función lineal:
\[
    f(x; \theta) = \mathrm{sign}(\theta_1 x_1 + \cdots + \theta_d x_d)
                 = \mathrm{sign}(\theta^{\mathsf{T}} x)
\]

Se toma el signo para nuestro problema porque solo estamos
considerando el caso binario: acepta o rechaza la entrada al edificio.

El modelo encontrará una línea, un plano que permita separar a las
personas que pueden o no entrar al edificio. Ese plano es justo el
lugar donde $\theta^{\mathsf{T}} x = 0$, y $\theta$ es el vector
normal a él, así que ajustar $\theta$ es girar y mover el plano.

El chiste es encontrar la función a partir de las entradas que tengo
en el entrenamiento. Cuando no es posible construir la función,
queremos entonces descubrir otra que cometa el menor error posible.

\begin{ojo}
  En nuestro problema el clasificador lineal no guarda info de los
  píxeles vecinos, solo evaluaría el píxel individual.
\end{ojo}

% =====================================================================

\subsection*{Funciones de pérdida}

Mencionamos que queremos encontrar el error. Esto se mide con la
diferencia del dato esperado con el dato obtenido, y hay diferentes
medidas de error.

A las funciones de error se les llama \textbf{loss}, función de error
o función de pérdida, y hacen lo que mencionábamos: dicen qué tanta
diferencia hay entre el dato observado y el esperado. Son esas
funciones las que se tratan de optimizar.

% =====================================================================

\subsection*{El perceptrón}

Queremos encontrar un clasificador lineal que minimice el error, o sea
encontrar $\theta$ que minimice el error de entrenamiento:
\[
    E(\theta) = \frac{1}{n} \sum_{i=1}^{n}
                \left[\, 1 - \delta\big(y^{(i)},\,
                f(x^{(i)}; \theta)\big) \right]
\]

Donde $\delta$ vale $1$ si sus dos argumentos coinciden y $0$ si no,
así que $1 - \delta$ vale $1$ exactamente cuando el ejemplo quedó mal
clasificado. Dicho de otra forma, $E(\theta)$ es la fracción de
ejemplos de entrenamiento en los que el modelo se equivoca.

¿Pero cómo construir ahora esa función lineal? ¿Cómo ajustar los
parámetros? ¿Hay alguna manera razonable de setear los parámetros de
$\theta$?

Una opción es ajustar incrementalmente cada parámetro para que veamos
que está bien. Un algoritmo simple de este tipo es el
\textbf{perceptron update rule}:
\[
    \text{si } y^{(i)} \big(\theta^{\mathsf{T}} x^{(i)}\big) \leq 0
    \quad\text{entonces}\quad
    \theta' = \theta + y^{(i)} x^{(i)}
\]

En otras palabras, los parámetros son actualizados solo si hay error, y
estos updates tienden a corregir errores. En nuestro caso, la
clasificación de caras, cuando el modelo hace un error el signo no
coincide con el de la etiqueta.

Supongamos que hacemos un error, y que los parámetros actualizados
están dados por $\theta'$. Si consideramos clasificar la misma imagen
después del update:

\begin{align*}
    y^{(i)} \big({\theta'}^{\mathsf{T}} x^{(i)}\big)
      &= y^{(i)} \big(
         (\theta + y^{(i)} x^{(i)})^{\mathsf{T}} x^{(i)} \big) \\
      &= y^{(i)} \big(\theta^{\mathsf{T}} x^{(i)}\big)
         + \big(y^{(i)}\big)^{2} \norm{x^{(i)}}^{2} \\
      &= y^{(i)} \big(\theta^{\mathsf{T}} x^{(i)}\big)
         + \norm{x^{(i)}}^{2}
\end{align*}

Como $y^{(i)} \in \set{-1, 1}$ su cuadrado siempre es $1$, y
$\norm{x^{(i)}}^{2} > 0$. O sea que el update le suma algo positivo,
empujando ese ejemplo hacia el lado correcto del plano. Eso es lo que
quiere decir que los updates tienden a corregir errores.

\begin{ojo}
  Esto da origen al reconocedor del perceptrón, y este reconocedor es
  útil para problemas lineales: solo puede separar con una línea o un
  plano, pero jamás podrá separar datos que no sean linealmente
  separables.
\end{ojo}

\newpage

% =====================================================================
%  Clase 01 Sep - Convergencia y generalizacion del perceptron
% =====================================================================

\clase{Convergencia y generalización (perceptrón)}{1 de septiembre
de 2026}

\begin{quick_sum}
    En esta sesion vimos la convergencia del perceptron, y la definicion
    formal, q era el margen geometrico y cerramos con la intro a SVM
\end{quick_sum}

\subsection*{Punto de partida y supuestos}

Seguimos con clasificadores lineales que pasan por el origen:
\[
    f(x; \theta) = \mathrm{sign}(\theta^{\mathsf{T}} x) ,
\]
donde $\theta$ denota los parámetros que tenemos que estimar en base a
los ejemplos de entrenamiento, con ejemplos $x$ y etiquetas $y$

Sea $k$ el número de parámetros actualizados y $\theta^{(k)}$ el
vector de parámetros después de $k$ actualizaciones. Cada vez que el
modelo se equivoca hace un update con la regla del perceptrón, y ese
update es un $k$ más

Para poder demostrar algo hacen falta dos supuestos:

\begin{enumerate}[itemsep=0.2em]
    \item \textbf{Las normas están acotadas}: existe $R$ tal que
          $\norm{x^{(t)}} \leq R$ para todo ejemplo.
    \item \textbf{El clasificador lineal existe}: hay un
          $\theta^{*}$, el $\theta$ ideal, que separa los datos con
          margen. Formalmente,
          \[
              y^{(t)} \big(\theta^{*}\big)^{\mathsf{T}} x^{(t)}
              \geq \gamma > 0
              \qquad \text{para todo } t ,
          \]
          donde $\gamma$ es el margen del clasificador, que es lo que
          nos dice que lo hace bien
\end{enumerate}

\begin{ojo}
  En el ejemplo de las imágenes de las caras reconocidas para entrar
  al edificio, el segundo supuesto es justamente suponer que ese
  clasificador existe. Si los datos no son linealmente separables, no
  hay $\theta^{*}$ y nada de lo que sigue aplica
\end{ojo}

% =====================================================================

\subsection*{Convergencia en un número finito de updates}

La convergencia se demuestra combinando dos resultados. El primero
dice que el producto interno con el ideal crece linealmente, y el
segundo que la norma cuadrada crece a lo más linealmente. Tomamos
$\theta^{(0)} = 0$.

\subsubsection*{Parte 1: el producto interno crece linealmente}

Supongamos que el update $k$ se hizo sobre el ejemplo
$(x^{(t)}, y^{(t)})$, o sea
$\theta^{(k)} = \theta^{(k-1)} + y^{(t)} x^{(t)}$. Entonces
\begin{align*}
    \big(\theta^{*}\big)^{\mathsf{T}} \theta^{(k)}
      &= \big(\theta^{*}\big)^{\mathsf{T}} \theta^{(k-1)}
         + y^{(t)} \big(\theta^{*}\big)^{\mathsf{T}} x^{(t)} \\
      &\geq \big(\theta^{*}\big)^{\mathsf{T}} \theta^{(k-1)}
            + \gamma ,
\end{align*}
donde la desigualdad es el supuesto del margen. Cada update le suma
al menos $\gamma$, así que aplicándolo $k$ veces desde
$\theta^{(0)} = 0$:
\[
    \big(\theta^{*}\big)^{\mathsf{T}} \theta^{(k)} \geq k\,\gamma .
\]

\subsubsection*{Parte 2: la norma cuadrada crece a lo más linealmente}

Con el mismo update,
\begin{align*}
    \norm{\theta^{(k)}}^{2}
      &= \norm{\theta^{(k-1)} + y^{(t)} x^{(t)}}^{2} \\
      &= \norm{\theta^{(k-1)}}^{2}
         + 2\,y^{(t)} \big(\theta^{(k-1)}\big)^{\mathsf{T}} x^{(t)}
         + \big(y^{(t)}\big)^{2} \norm{x^{(t)}}^{2} \\
      &\leq \norm{\theta^{(k-1)}}^{2} + R^{2} .
\end{align*}

El término de en medio se va porque el update ocurre precisamente
cuando el ejemplo estaba mal clasificado, o sea cuando
$y^{(t)} \big(\theta^{(k-1)}\big)^{\mathsf{T}} x^{(t)} \leq 0$. Y el
último se acota porque $\big(y^{(t)}\big)^{2} = 1$ y
$\norm{x^{(t)}} \leq R$. Repitiendo $k$ veces:
\[
    \norm{\theta^{(k)}}^{2} \leq k\,R^{2} 
\]

\subsubsection*{Combinación: qué pasa con el ángulo}

Si combinamos esos dos resultados podemos ver el coseno del ángulo
entre el vector de la actualización $k$ y el ideal:
\[
    \cos\!\big(\theta^{*}, \theta^{(k)}\big) =
    \frac{\big(\theta^{*}\big)^{\mathsf{T}} \theta^{(k)}}
         {\norm{\theta^{*}}\,\norm{\theta^{(k)}}}
    \geq \frac{k\,\gamma}{\norm{\theta^{*}}\,\sqrt{k}\,R}
    = \sqrt{k}\;\frac{\gamma}{\norm{\theta^{*}}\,R} .
\]

El numerador se acotó por abajo con la parte 1 y el denominador por
arriba con la parte 2. Como el coseno está acotado por $1$:
\[
    \sqrt{k}\;\frac{\gamma}{\norm{\theta^{*}}\,R} \leq 1
    \qquad\Longrightarrow\qquad
    k \leq \frac{R^{2}\,\norm{\theta^{*}}^{2}}{\gamma^{2}} .
\]

\begin{idea}
  Cada update acerca el coseno a $1$ en una cantidad finita, o sea
  acerca el vector actual al ideal. Como el coseno no puede pasar de
  $1$, solo se puede hacer un número finito de updates. Ahí está la
  convergencia
\end{idea}

\begin{ojo}
  Está acotado, pero la cota depende del valor de los pesos: puede
  ser muy grande o muy pequeña, y de antemano no sabemos cuál.
\end{ojo}

% =====================================================================

\subsection*{Margen geométrico}

El cociente $\norm{\theta^{*}}^{2} / \gamma^{2}$ nos dice qué tan
difícil es el problema, o sea cuántos updates voy a tener que hacer.
Su inverso es la distancia más chica que hay entre una instancia del
conjunto de entrenamiento y mi región de decisión

Si esa distancia es grande hay bastante separación entre la
superficie que divide las clases y las instancias, y es más fácil
clasificar. Si es chica, es más difícil

\subsubsection*{Cómo se calcula}

Se mide la distancia de la superficie de decisión a uno de los datos
de entrenamiento. Sea $\tilde{x}$ el punto de la frontera más cercano
a $x^{(t)}$, y $\varepsilon$ esa distancia. Moverse de $x^{(t)}$ a la
frontera es avanzar $\varepsilon$ en la dirección normal, que es la
de $\theta^{*}$:
\[
    \tilde{x} = x^{(t)} - y^{(t)}\,\varepsilon\,
                \frac{\theta^{*}}{\norm{\theta^{*}}} .
\]

Evaluando el clasificador ideal en ese punto:
\[
    y^{(t)} \big(\theta^{*}\big)^{\mathsf{T}} \tilde{x}
    = \underbrace{y^{(t)}
      \big(\theta^{*}\big)^{\mathsf{T}} x^{(t)}}_{=\;\gamma}
      - \varepsilon\,\norm{\theta^{*}}
    = \gamma - \varepsilon\,\norm{\theta^{*}} = 0 ,
\]
porque $\tilde{x}$ está sobre la frontera y ahí el clasificador vale
cero. Despejando queda el \textbf{margen geométrico}:
\[
    \varepsilon = \frac{\gamma}{\norm{\theta^{*}}} .
\]

\begin{figure}[htbp]
  \centering
  \begin{tikzpicture}[>=stealth]
    % Frontera de decision
    \draw[acento, very thick] (-1.6,-1.6) -- (2.7,2.7);
    \node[acento, font=\footnotesize, anchor=south east]
      at (2.7,2.75) {$(\theta^{*})^{\mathsf{T}} x = 0$};

    % Vector normal
    \draw[->, red!70!black, thick] (0,0) -- (0.9,-0.9);
    \node[red!70!black, font=\footnotesize, anchor=north west]
      at (0.85,-0.85) {$\theta^{*}$};

    % Punto de entrenamiento y su proyeccion
    \coordinate (xt) at (2.6,-0.4);
    \coordinate (xp) at (1.1,1.1);
    \draw[dashed, gray!60!black] (xt) -- (xp);
    \fill[red!70!black] (xt) circle (2.2pt);
    \fill[acento] (xp) circle (1.8pt);
    \node[font=\footnotesize, anchor=west] at (2.75,-0.4)
      {$x^{(t)}$};
    \node[font=\footnotesize, anchor=south east] at (1.05,1.15)
      {$\tilde{x}$};
    \node[font=\footnotesize, anchor=south west] at (1.9,0.25)
      {$\varepsilon$};

    % Ejemplos de las dos clases
    \fill[red!70!black] (2.0,-1.2) circle (1.6pt);
    \fill[red!70!black] (2.9,0.6) circle (1.6pt);
    \fill[acento] (-0.6,1.2) circle (1.6pt);
    \fill[acento] (0.4,2.3) circle (1.6pt);
  \end{tikzpicture}
\end{figure}

Por lo tanto el $k$, el número de veces que tengo que actualizar, se
puede expresar en términos del margen geométrico:
\[
    k \leq \frac{R^{2}}{\varepsilon^{2}} .
\]

\begin{idea}
  Entre más ancho el pasillo que separa las clases, menos updates
  necesita el perceptrón. La dificultad del problema no está en
  cuántos datos hay ni en cuántas dimensiones, sino en qué tan
  apretado quedó ese pasillo
\end{idea}

% =====================================================================

\subsection*{Generalización}

Hasta ahorita asumimos que el conjunto de datos era linealmente
separable, pero qué pasa si ahora no trabajo con mi conjunto de
entrenamiento sino con datos que están fuera de él

Cuando llegue una instancia nueva, el clasificador lineal tendrá que
ser capaz de clasificarla: se pone en el plano y se ve si la
superficie separadora la separa correctamente


\subsubsection*{La garantía formal}

El argumento se puede hacer preciso. Supongamos que \emph{todas} las
imágenes y etiquetas que podríamos llegar a encontrar, no solo las del
entrenamiento, cumplen los mismos dos supuestos: $\norm{x_t} < R$, y
existe un $\theta^{*}$ finito con
$y_t (\theta^{*})^{\mathsf{T}} x_t \geq \gamma$ para toda $t$. O sea
que existe un clasificador lineal que sirve para todo el problema, no
solo para la muestra.

Ahora imaginemos que las instancias llegan una por una, y que por cada
una hacemos a lo más un update si la clasificamos mal, y seguimos
adelante sin volver a verla. ¿Cuántos errores vamos a cometer en esa
secuencia arbitraria e indefinida?

\begin{idea}
  El mismo número: $k \leq (R / \gamma_{\mathrm{geom}})^{2}$. La cota
  no cambia. Una vez cometidos esos errores, todas las instancias
  nuevas se clasifican correctamente. Eso es una garantía de
  generalización, y no salió de un supuesto estadístico sobre la
  distribución de los datos, sino de la geometría del problema.
\end{idea}

\begin{ojo}
  Con un pero: el algoritmo necesita \emph{saber} cuándo se equivocó,
  porque la cota está expresada en número de updates, y los updates se
  hacen sobre errores. En un escenario real donde nadie te dice la
  etiqueta verdadera, no hay quién dispare el update.
\end{ojo}

\begin{idea}
  Lo que queremos ahora es un clasificador lineal que nos dé el
  máximo margen geométrico. De todas las superficies que separan bien
  el entrenamiento, la que deja el pasillo más ancho es la que tiene
  más holgura para los datos nuevos. Ese clasificador son las
  máquinas de vectores de soporte
\end{idea}

\newpage

% =====================================================================
%  Clase 03 Sep - SVM
% =====================================================================

\clase{Máquinas de vectores de soporte (SVM)}{3 de septiembre de 2026}

\begin{quick_sum}
  En esta clase vimos como funcionan los SVM
\end{quick_sum}

\subsection*{El problema}

Se plantea el problema de encontrar la superficie de separación entre
los datos de una muestra, de tal manera que el margen geométrico que
vimos sea el máximo posible

\begin{figure}[htbp]
  \centering
  \begin{tikzpicture}[>=stealth]
    % Frontera y margenes
    \draw[acento, very thick] (0,0) -- (3.3,3.3);
    \node[acento, font=\footnotesize, anchor=west] at (3.35,3.3)
      {$\theta^{\mathsf{T}} x + \theta_0 = 0$};
    \draw[gray!70!black, dashed] (0,1) -- (2.5,3.5);
    \draw[gray!70!black, dashed] (1,0) -- (3.5,2.5);

    % Margen
    \draw[<->, red!70!black, thick] (1,2) -- (2,1);
    \node[red!70!black, font=\footnotesize, anchor=west]
      at (1.62,1.75) {$2\varepsilon$};

    % Clase -1, arriba
    \fill[acento] (1,2) circle (2.4pt);
    \draw[acento, thick] (1,2) circle (4.5pt);
    \fill[acento] (0.4,1.4) circle (2.4pt);
    \draw[acento, thick] (0.4,1.4) circle (4.5pt);
    \fill[acento] (0.3,2.6) circle (2.2pt);
    \fill[acento] (1.3,3.1) circle (2.2pt);

    % Clase +1, abajo
    \fill[red!70!black] (2,1) circle (2.4pt);
    \draw[red!70!black, thick] (2,1) circle (4.5pt);
    \fill[red!70!black] (2.9,1.9) circle (2.4pt);
    \draw[red!70!black, thick] (2.9,1.9) circle (4.5pt);
    \fill[red!70!black] (3.1,0.7) circle (2.2pt);
    \fill[red!70!black] (2.5,0.3) circle (2.2pt);
  \end{tikzpicture}
  \caption{La superficie que deja el pasillo más ancho. Los puntos
    circulados, los que caen justo sobre el borde del margen, son los
    vectores de soporte.}
  \label{fig:svm}
\end{figure}

Se plantea como un problema de optimización: lo que quiero optimizar
son los parámetros de mi vector, de tal manera que exista un margen
alrededor de mi superficie de separación que sea lo más grande
posible.

Una manera de lograrlo sería encontrar una superficie que separe bien
las instancias y luego tratar de que ese clasificador tenga el mayor
margen posible. Se busca entonces maximizar el margen geométrico
$\gamma / \norm{\theta}$, cumpliendo las restricciones
\[
    y^{(t)} \big(\theta^{\mathsf{T}} x^{(t)}\big) \geq \gamma
    \qquad \text{para } t = 1, 2, \ldots, n .
\]

\begin{idea}
  Esa restricción dice dos cosas al mismo tiempo. El \emph{signo} de
  $y^{(t)}\theta^{\mathsf{T}}x^{(t)}$ dice si la instancia quedó del
  lado correcto del plano, y su \emph{magnitud} dice qué tan lejos
  quedó. Pedir que sea $\geq \gamma$ es pedir las dos: bien
  clasificada, y además a distancia de sobra, no pegada a la
  frontera.
\end{idea}

% =====================================================================

\subsection*{Como problema de optimización}

Maximizar el cociente que dijimos es lo mismo que minimizar su
inverso al cuadrado:
\[
    \min_{\theta} \; \frac{1}{2}\,
    \frac{\norm{\theta}^{2}}{\gamma^{2}}
    \qquad \text{sujeto a} \qquad
    y^{(t)} \theta^{\mathsf{T}} x^{(t)} \geq \gamma ,
    \quad t = 1, \ldots, n .
\]

Así escrito el problema depende de $\gamma$, y eso sobra. Si se
escala $\theta$ y $\gamma$ por la misma constante, ni las
restricciones ni el margen geométrico cambian: la superficie es la
misma. Se aprovecha esa libertad para fijar $\gamma = 1$, y el
problema queda sin $\gamma$:
\[
    \min_{\theta} \; \frac{1}{2}\,\norm{\theta}^{2}
    \qquad \text{sujeto a} \qquad
    y^{(t)} \theta^{\mathsf{T}} x^{(t)} \geq 1 ,
    \quad t = 1, \ldots, n .
\]

Con esa normalización el margen geométrico es $1 / \norm{\theta}$,
así que minimizar la norma es literalmente ensanchar el pasillo.

\subsubsection*{El offset}

Hasta ahora nuestra función está atada a pasar por el origen, y no
queremos eso. Para lograrlo se suma un vector de offset $\theta_0$:
\[
    f(x; \theta, \theta_0) =
    \mathrm{sign}\big(\theta^{\mathsf{T}} x + \theta_0\big) ,
\]
y el problema se vuelve
\[
    \min_{\theta,\,\theta_0} \; \frac{1}{2}\,\norm{\theta}^{2}
    \qquad \text{sujeto a} \qquad
    y^{(t)} \big(\theta^{\mathsf{T}} x^{(t)} + \theta_0\big) \geq 1 .
\]

Esto ayuda a que podamos encontrar una superficie con mayor margen
geométrico, porque ya no obligamos al plano a pasar por el origen.

\subsubsection*{Vectores de soporte}

Los elementos que caen justo a la distancia geométrica son los
llamados \textbf{vectores de soporte}, y son elementos de la muestra.
Son los que están circulados en la figura \ref{fig:svm}.


\subsection*{Propiedades del clasificador de margen máximo}

\textbf{A favor:}

\begin{itemize}[itemsep=0.2em]
    \item \textbf{La solución es única} para cualquier conjunto de
          entrenamiento linealmente separable. El perceptrón, en
          cambio, se detiene en cualquier separador que funcione, y
          cuál depende del orden en que le llegaron los datos.
    \item \textbf{Es robusto a ejemplos ruidosos}, porque dejar la
          frontera lo más lejos posible de los datos da holgura para
          que una instancia corrida de lugar siga cayendo del lado
          correcto.
    \item \textbf{Depende solo de un subconjunto} de los ejemplos, los
          que caen exactamente sobre el margen. Los demás podrían
          estar en cualquier parte fuera del margen sin cambiar la
          solución.
\end{itemize}

\textbf{En contra:}

\begin{ojo}
  Es robusto a ejemplos ruidosos, pero \textbf{no a etiquetas
  ruidosas}. Un solo ejemplo mal etiquetado puede cambiar
  radicalmente el clasificador de margen máximo: si queda del lado
  equivocado, el pasillo tiene que encogerse hasta acomodarlo, y la
  frontera se va a otro lado. Esa fragilidad es la que motiva las
  variables de holgura.
\end{ojo}

% =====================================================================

\subsection*{Validación cruzada}

Para evaluar este clasificador podemos hacer uso de la validación
cruzada. Consiste en tomar el conjunto de entrenamiento, separar un
cachito, construir el modelo con el resto y evaluar en ese cachito.
El cachito se llama \textbf{fold}, y esto se repite $k$ veces
cambiando cuál fold se deja fuera.

\begin{figure}[htbp]
  \centering
  \begin{tikzpicture}
    \foreach \i in {1,...,5} {
      \node[anchor=east, font=\footnotesize] at (0,{-0.62*\i})
        {Iteración \i};
      \foreach \j in {1,...,5} {
        \ifnum\i=\j
          \draw[fill=fondoOjo, draw=orange!70!black]
            ({0.25+1.2*(\j-1)},{-0.62*\i-0.2})
            rectangle ++(1.1,0.4);
        \else
          \draw[fill=fondoIdea, draw=acento!60]
            ({0.25+1.2*(\j-1)},{-0.62*\i-0.2})
            rectangle ++(1.1,0.4);
        \fi
      }
    }
    \draw[fill=fondoOjo, draw=orange!70!black]
      (0.25,-4.2) rectangle ++(0.45,0.3);
    \node[anchor=west, font=\footnotesize] at (0.8,-4.05) {prueba};
    \draw[fill=fondoIdea, draw=acento!60]
      (2.4,-4.2) rectangle ++(0.45,0.3);
    \node[anchor=west, font=\footnotesize] at (2.95,-4.05)
      {entrenamiento};
  \end{tikzpicture}
  \caption{Validación cruzada con $k = 5$. En cada iteración un fold
    distinto hace de conjunto de prueba.}
  \label{fig:kfolds}
\end{figure}

\subsubsection*{Leave-one-out}

Si dejamos solo un elemento fuera es el \textbf{leave-one-out cross
validation}, pero eso se hace $n$ veces, con $n$ el número de
elementos del training set. El error se mide así:
\[
    E_{\mathrm{loo}} = \frac{1}{n} \sum_{i=1}^{n}
    \left[\, 1 - \delta\big(y^{(i)},\,
    f(x^{(i)};\, \theta^{(-i)})\big) \right] ,
\]
donde $\theta^{(-i)}$ son los parámetros entrenados con todo el
conjunto \emph{menos} el ejemplo $i$. O sea que cada ejemplo se
predice con un modelo que nunca lo vio, y $E_{\mathrm{loo}}$ es la
fracción de veces que ese modelo se equivoca.

\begin{idea}
  El leave-one-out error está acotado por el número de vectores de
  soporte:
  \[
      E_{\mathrm{loo}} \leq \frac{\#\text{vectores de soporte}}{n} .
  \]
  Cada vez que se quita una instancia que sea vector de soporte se
  puede cometer un error, porque quitarla mueve la superficie. Pero
  quitar una que no lo sea no cambia nada, así que ese ejemplo se
  sigue clasificando bien. Un clasificador con pocos vectores de
  soporte generaliza mejor.
\end{idea}

% =====================================================================

\subsection*{Datos con ruido: variables de holgura}

Uno de los problemas de este clasificador es que si las etiquetas del
conjunto están mal, el clasificador va a estar incorrecto. También
hay casos en los que el dato mismo trae error.

Para solucionar esto se introducen \textbf{variables de holgura},
dando una optimización relajada:
\[
    \min_{\theta,\,\theta_0,\,\xi} \; \frac{1}{2}\,\norm{\theta}^{2}
    + C \sum_{t=1}^{n} \xi_t
\]
sujeta a
\[
    y^{(t)} \big(\theta^{\mathsf{T}} x^{(t)} + \theta_0\big)
    \geq 1 - \xi_t ,
    \qquad \xi_t \geq 0 ,
    \qquad t = 1, \ldots, n .
\]

Cada $\xi_t$ mide cuánto se le perdona a la instancia $t$, y $C$ es
un parámetro de costo, o sea qué tanto le va a costar violar las
restricciones.

\begin{ojo}
  Si el costo no es tan grande, le voy a dar mucho chance a que
  queden instancias dentro del margen geométrico. Y al revés: con $C$
  grande no se permiten violaciones, el modelo se aprieta contra el
  entrenamiento y vuelve a ser sensible al ruido.
\end{ojo}

% =====================================================================

\subsection*{Cuando los datos no son linealmente separables}

Todo lo anterior es solo para instancias linealmente separables, pero
qué pasa si las instancias no lo son. A alguien se le ocurrió llevar
las instancias que no son linealmente separables a un espacio en el
que sí lo sean.

Entonces, en lugar de trabajar con la instancia del vector original,
se trabaja con su imagen $\varphi(x)$ en el nuevo espacio. Ahí se
construye el clasificador lineal y ahí se obtiene la respuesta:
\[
    y^{(t)} \big(\theta^{\mathsf{T}} \varphi(x^{(t)}) + \theta_0\big)
    \geq 1 - \xi_t .
\]

\begin{idea}
  La superficie sigue siendo un plano, pero en el espacio nuevo. Al
  regresarla al espacio original se ve curva, y por eso el modelo
  puede separar datos que una recta no separaba.
\end{idea}

\newpage

% =====================================================================
%  Clase 08 Sep - SVM revisitada: forma dual, kernels y multiclase
% =====================================================================

\clase{SVM revisitada: forma dual y kernels}{8 de septiembre de 2026}

\begin{quick_sum}
    Retomamos el problema de mapear las instancias a un espacio donde
    sí sean separables, y vimos que el mapeo explícito es caro y
    normalmente ni siquiera lo conocemos. La salida es reescribir la
    SVM en su forma dual, donde los ejemplos aparecen solo dentro de
    productos internos, y ahí sustituir el producto por un kernel.
    Cerramos con cómo se extiende todo esto a más de dos clases.
\end{quick_sum}

\subsection*{El mapeo explícito y su costo}

La idea de la clase pasada: si el conjunto no es linealmente separable
en $\R^d$, se busca una función $\varphi$ que lleve las instancias a un
espacio de mayor dimensión donde sí lo sean,
\[
    \varphi : \R^{d} \to \R^{\mathcal{F}} ,
    \qquad \mathcal{F} > d .
\]

\begin{ejemplo}
  De dos dimensiones a tres, elevando y multiplicando coordenadas:
  \[
      [x_1, x_2] \;\longmapsto\;
      \varphi([x_1, x_2]) = [\,x_1^{2},\; x_2^{2},\; x_1 x_2\,] .
  \]
  El producto interno en el espacio nuevo queda
  \[
      \varphi(x)^{\mathsf{T}} \varphi(z) =
      x_1^{2} z_1^{2} + x_2^{2} z_2^{2} + x_1 x_2 z_1 z_2 .
  \]
  Con tres dimensiones de entrada el mismo esquema da nueve
  coordenadas, todos los productos $x_i x_j$.
\end{ejemplo}

\begin{ojo}
  Aquí está el problema. Encontrar la $\varphi$ adecuada no es
  trivial, y aunque se encuentre, calcular productos internos en
  $\R^{\mathcal{F}}$ puede salir carísimo: la dimensión crece rápido.
  Peor aún, en general \emph{no conocemos} $\varphi$ explícitamente.
\end{ojo}

La salida no es buscar mejores mapeos, sino reformular la SVM de modo
que $\varphi$ nunca aparezca sola, solo dentro de productos internos.
Eso es la forma dual.

% =====================================================================

\subsection*{La forma dual}

El problema en el espacio de características es el de siempre:
\[
    \min_{\theta,\,\theta_0} \; \frac{\norm{\theta}^{2}}{2}
    \qquad \text{sujeto a} \qquad
    y_t\big(\theta^{\mathsf{T}} \varphi(x_t) + \theta_0\big) \geq 1 ,
    \quad t = 1, \ldots, n .
\]

Es convexo y con restricciones lineales, y ese tipo de problemas se
puede pasar a su forma dual con multiplicadores de Lagrange. Se
introduce un escalar no negativo $\alpha_t$ por cada restricción:
\[
    J(\theta, \theta_0; \alpha) = \frac{\norm{\theta}^{2}}{2}
    - \sum_{t=1}^{n} \alpha_t
      \Big[\, y_t\big(\theta^{\mathsf{T}} \varphi(x_t)
      + \theta_0\big) - 1 \,\Big] .
\]

\begin{idea}
  Cada $\alpha_t$ es el precio de violar la restricción $t$. Si la
  restricción se cumple con holgura, al maximizar sobre $\alpha$ le
  conviene poner $\alpha_t = 0$: esa restricción no está apretando y
  no cobra nada. Solo las que se cumplen con igualdad, o sea las que
  tocan el margen, terminan con $\alpha_t > 0$.
\end{idea}

Bajo las condiciones de optimalidad de Lagrange se pueden intercambiar
el mínimo y el máximo sin cambiar la respuesta:
\[
    \underbrace{\min_{\theta,\theta_0} \; \max_{\alpha \geq 0}
    J(\theta, \theta_0; \alpha)}_{\text{forma primal}}
    \;=\;
    \underbrace{\max_{\alpha \geq 0} \; \min_{\theta,\theta_0}
    J(\theta, \theta_0; \alpha)}_{\text{forma dual}} .
\]

Resolviendo primero el mínimo interno, o sea derivando e igualando a
cero:
\begin{align*}
    \frac{\partial J}{\partial \theta_0}
      &= - \sum_{t=1}^{n} \alpha_t y_t = 0 , \\
    \frac{\partial J}{\partial \theta}
      &= \theta - \sum_{t=1}^{n} \alpha_t y_t \varphi(x_t) = 0
      \;\Longrightarrow\;
      \theta = \sum_{t=1}^{n} \alpha_t y_t \varphi(x_t) .
\end{align*}

\begin{idea}
  Esa segunda ecuación dice algo fuerte: la solución $\theta$ siempre
  vive en el subespacio generado por los ejemplos de entrenamiento. No
  es un vector cualquiera de $\R^{\mathcal{F}}$, es una combinación
  lineal de las propias instancias, pesada por $\alpha_t y_t$.
\end{idea}

Sustituyendo esa forma de $\theta$ en el objetivo se llega al dual:
\[
    \max_{\alpha} \;\; \sum_{t=1}^{n} \alpha_t
    - \frac{1}{2} \sum_{i=1}^{n} \sum_{j=1}^{n}
      \alpha_i \alpha_j y_i y_j
      \big[\varphi(x_i)^{\mathsf{T}} \varphi(x_j)\big]
\]
sujeto a
\[
    \alpha_t \geq 0 , \qquad \sum_{t=1}^{n} \alpha_t y_t = 0 .
\]

Sigue siendo un problema de programación cuadrática, pero las
restricciones son mucho más simples que las del primal.

\subsubsection*{La matriz de Gram}

Fíjate en lo que quedó: la dimensión $\mathcal{F}$ del espacio de
características \textbf{ya no aparece}. El problema está formulado
enteramente sobre la matriz de productos internos, llamada matriz de
Gram:
\[
    K =
    \begin{bmatrix}
      \varphi(x_1)^{\mathsf{T}}\varphi(x_1) & \cdots &
      \varphi(x_1)^{\mathsf{T}}\varphi(x_n) \\
      \vdots & \ddots & \vdots \\
      \varphi(x_n)^{\mathsf{T}}\varphi(x_1) & \cdots &
      \varphi(x_n)^{\mathsf{T}}\varphi(x_n)
    \end{bmatrix} .
\]

Es de $n \times n$, o sea que su tamaño depende del número de
ejemplos, no de la dimensión del espacio. Ahí está la puerta.

% =====================================================================

\subsection*{El truco del kernel}

\begin{idea}
  \textbf{El truco del kernel.} Si tienes un algoritmo donde los
  ejemplos aparecen \emph{solo} dentro de productos internos, puedes
  reemplazar ese producto interno por otro, sin siquiera conocer
  $\varphi$.
\end{idea}

Se define entonces una función kernel $K(x, x')$ tal que
\[
    K(x, x') = \varphi(x)^{\mathsf{T}} \varphi(x') ,
\]
y se sustituye. Nunca se calcula $\varphi(x)$: basta con conocer la
forma funcional del kernel. Lo que costaba $\mathcal{F}$ operaciones
ahora cuesta lo que cueste evaluar $K$, que suele ser $\Ord{d}$.

\begin{teorema}[Mercer]
  Para toda función $K : \R^{d} \times \R^{d} \to \R$ simétrica y
  semidefinida positiva, existen un espacio de Hilbert $\mathcal{F}$ y
  una función $\varphi : \R^{d} \to \mathcal{F}$ tales que
  \[
      K(x, x') = \langle \varphi(x), \varphi(x') \rangle .
  \]
\end{teorema}

\begin{idea}
  Este teorema es el permiso para trabajar al revés. No hace falta
  construir el kernel a partir de un conjunto de funciones base
  $\varphi_1(x), \ldots, \varphi_r(x)$: basta definir una $K$ que
  cumpla las dos condiciones, simetría y semidefinida positiva, y el
  teorema garantiza que \emph{alguna} $\varphi$ existe. Nunca hay que
  verla.
\end{idea}

\subsubsection*{Kernels estándar}

\begin{center}
\begin{tabular}{ll}
    \toprule
    \textbf{Polinomial} &
      $K(x, x') = \big(1 + x^{\mathsf{T}} x'\big)^{p}$,
      \quad $p = 1, 2, \ldots$ \\[0.4em]
    \textbf{Base radial (RBF)} &
      $K(x, x') = \exp\!\left(
      -\dfrac{\beta}{2}\,\norm{x - x'}^{2}\right)$,
      \quad $\beta > 0$ \\
    \bottomrule
\end{tabular}
\end{center}

Con $p = 1$ el polinomial es prácticamente el producto interno
original, o sea la SVM lineal de siempre. El RBF corresponde a un
espacio de características de dimensión infinita, y aun así se evalúa
con una resta, una norma y una exponencial.

% =====================================================================

\subsection*{Clasificar con la solución dual}

Ya con los $\alpha_t$ resueltos, una instancia nueva se clasifica con
la función discriminante:
\[
    \hat{y}(x) = \theta^{\mathsf{T}}\varphi(x) + \theta_0
    = \sum_{t \in \mathrm{SV}} \alpha_t y_t\, K(x_t, x) + \theta_0 ,
\]
donde $\mathrm{SV}$ es el conjunto de vectores de soporte, o sea los
que tienen $\alpha_t \neq 0$.

\begin{idea}
  Aquí reaparece la dispersión que ya conocíamos. Muchos $\alpha_t$
  salen exactamente cero de la optimización, y son distintos de cero
  \emph{solo} para los vectores de soporte. Por eso la suma corre
  sobre $\mathrm{SV}$ y no sobre los $n$ ejemplos: los demás no
  aportan nada al clasificar.
\end{idea}

\begin{ojo}
  No sabemos qué ejemplos van a ser vectores de soporte hasta haber
  resuelto la optimización. Y su identidad depende del kernel: cambiar
  de kernel cambia quiénes son.
\end{ojo}

\subsubsection*{De dónde sale $\theta_0$}

El offset se cayó del problema dual, así que hay que recuperarlo
aparte. Se usa que para todo vector de soporte la restricción se
cumple con igualdad:
\[
    y_i\Big(\sum_{t \in \mathrm{SV}} \alpha_t y_t\, K(x_t, x_i)
    + \theta_0\Big) = 1 ,
    \qquad i \in \mathrm{SV} ,
\]
y como $y_i \in \set{-1, 1}$ implica $y_i^{2} = 1$, se despeja
\[
    \theta_0 = y_i - \sum_{t \in \mathrm{SV}}
    \alpha_t y_t\, K(x_t, x_i) .
\]

\subsubsection*{El margen geométrico en forma kernel}

Sigue siendo $1/\norm{\theta}$, solo que ahora $\norm{\theta}$ se
escribe con el kernel:
\[
    \hat{\gamma}_{\mathrm{geom}} = \left(
    \sum_{i=1}^{n} \sum_{j=1}^{n}
    \alpha_i \alpha_j y_i y_j\, K(x_i, x_j)
    \right)^{-1/2} .
\]

% =====================================================================

\subsection*{La forma dual relajada}

Todo lo anterior supone que los datos sí son separables en el espacio
de características. Con variables de holgura el primal era
\[
    \min \; \frac{\norm{\theta}^{2}}{2} + C \sum_{t=1}^{n} \xi_t
    \qquad \text{sujeto a} \qquad
    y_t\big(\theta^{\mathsf{T}}\varphi(x_t) + \theta_0\big)
    \geq 1 - \xi_t .
\]

Su dual es casi idéntico al de arriba. La \emph{única} diferencia es
que los multiplicadores quedan acotados por arriba:
\[
    \max_{\alpha} \;\; \sum_{t=1}^{n} \alpha_t
    - \frac{1}{2} \sum_{i=1}^{n} \sum_{j=1}^{n}
      \alpha_i \alpha_j y_i y_j\, K(x_i, x_j)
    \qquad \text{sujeto a} \qquad
    0 \leq \alpha_t \leq C ,
    \quad \sum_{t=1}^{n} \alpha_t y_t = 0 .
\]

\begin{idea}
  Por qué la cota superior: sin ella, una restricción imposible de
  satisfacer haría crecer su $\alpha_t$ hasta infinito, porque el
  precio de violarla nunca dejaría de subir. El valor $C$ es
  justamente el punto donde uno decide dejar de insistir con esa
  restricción. Es el mismo $C$ del primal, visto desde el otro lado.
\end{idea}

\begin{ojo}
  Con holgura, $\theta_0$ ya no se despeja con cualquier $\alpha_t > 0$.
  Hay que usar los que cumplen $0 < \alpha_t < C$, o sea excluir los
  que están en el tope. Un $\alpha_t = C$ marca un punto que
  \emph{viola} la restricción de margen, y para ese la igualdad no se
  cumple.
\end{ojo}

\begin{center}
\begin{tabular}{lll}
    \toprule
    \textbf{Valor de $\alpha_t$} & \textbf{Dónde cae el ejemplo} &
      \textbf{Sirve para $\theta_0$} \\
    \midrule
    $\alpha_t = 0$     & Fuera del margen, bien clasificado & No \\
    $0 < \alpha_t < C$ & Justo sobre el margen              & Sí \\
    $\alpha_t = C$     & Viola el margen                    & No \\
    \bottomrule
\end{tabular}
\end{center}

% =====================================================================

\subsection*{Clasificación multiclase}

Hasta aquí todo fue binario. En multiclase cada punto de entrenamiento
pertenece a una de $N$ clases distintas, y se quiere una función que,
dado un punto nuevo, determine correctamente a cuál pertenece.

\begin{ojo}
  Eso \textbf{no} es lo mismo que multietiqueta. Hay escenarios con
  varias categorías donde un punto puede pertenecer a varias a la vez,
  y ese es otro problema.
\end{ojo}

\subsubsection*{One vs All (OVA)}

Se construyen $N$ clasificadores binarios. Para el $i$-ésimo, los
positivos son todos los puntos de la clase $i$ y los negativos todos
los que no. Si $f_i$ es ese clasificador, se decide con
\[
    f(x) = \arg\max_{i} \; f_i(x) .
\]

\subsubsection*{All vs All (AVA)}

Un clasificador por cada par de clases $i, j$, donde $i$ son los
positivos y $j$ los negativos. Como $f_{ji} = -f_{ij}$, en realidad
solo hay que entrenar $N(N-1)/2$ de ellos. Se decide sumando votos:
\[
    f(x) = \arg\max_{i} \; \sum_{j} f_{ij}(x) .
\]
También se le llama \emph{all-pairs} o \emph{one-vs-one}.

\begin{idea}
  OVA y AVA son tan simples que mucha gente los inventó por separado.
  Y eligiendo buenos clasificadores binarios de base, funcionan tan
  bien como cualquier otra cosa que se te ocurra.
\end{idea}

\begin{center}
\begin{tabular}{lll}
    \toprule
     & \textbf{OVA} & \textbf{AVA} \\
    \midrule
    Clasificadores & $\Ord{N}$ & $\Ord{N^{2}}$ \\
    Tamaño de cada uno & Todo el conjunto & Solo dos clases \\
    Balance de clases & Desbalanceado & Balanceado \\
    \bottomrule
\end{tabular}
\end{center}

AVA sale más rápido y usa menos memoria pese a entrenar más
clasificadores, porque cada uno es mucho más chico. Si el tiempo de
entrenamiento es superlineal en el número de puntos, que es el caso de
la SVM, AVA es la mejor opción. OVA además produce muestras de
entrenamiento desbalanceadas: una clase contra las otras $N-1$ juntas.

\subsubsection*{SVM multiclase directa}

En lugar de armar varios clasificadores binarios, lo natural es
distinguir todas las clases en una sola optimización.

\textbf{Weston y Watkins (1999).} Un vector de pesos $w_m$ por clase, y
una holgura $\xi_{i,t}$ por cada ejemplo y cada clase equivocada:
\[
    \min \;\; \frac{1}{2}\sum_{m=1}^{k} w_m^{\mathsf{T}} w_m
    + C \sum_{i=1}^{l} \sum_{t \neq y_i} \xi_{i,t}
\]
sujeto a
\[
    w_{y_i}^{\mathsf{T}} \varphi(x_i) + b_{y_i} \geq
    w_{t}^{\mathsf{T}} \varphi(x_i) + b_{t} + 2 - \xi_{i,t} ,
    \qquad \xi_{i,t} \geq 0 .
\]
Se decide con $\arg\max_m \big(w_m^{\mathsf{T}}\varphi(x) + b_m\big)$.
Su desventaja principal es el tiempo de cómputo: el problema
cuadrático resultante es enorme.

\textbf{Crammer y Singer (2001).} Misma idea pero con una sola holgura
$\xi_i$ por ejemplo, no una por clase:
\[
    \min \;\; \frac{1}{2}\sum_{m=1}^{k} w_m^{\mathsf{T}} w_m
    + C \sum_{i=1}^{l} \xi_i
\]
sujeto a
\[
    w_{y_i}^{\mathsf{T}} \varphi(x_i)
    - w_{t}^{\mathsf{T}} \varphi(x_i)
    \geq 1 - \delta_{y_i, t} - \xi_i ,
\]
donde $\delta_{i,j}$ es la delta de Kronecker, que vale $1$ si $i = j$
y $0$ si no. La optimización se hace sin los términos de offset $b_m$.

\textbf{SimMSVM (He et al., 2012).} Reduce el número de restricciones
de margen a $l$, una por ejemplo, comparando la clase correcta contra
el promedio de las demás:
\[
    w_{y_i}^{\mathsf{T}} \varphi(x_i)
    - \frac{1}{k-1} \sum_{m \neq y_i}
      w_m^{\mathsf{T}} \varphi(x_i) \geq 1 - \xi_i .
\]
Sus autores probaron que su formulación dual es exactamente la de
Crammer y Singer con una restricción adicional, y cada modelo de clase
se construye sobre sus propios vectores de soporte.

\newpage

% =====================================================================
%  Clase - Preparacion de datos y redes neuronales
% =====================================================================

\clase{Preparación de datos y redes neuronales}{10 de septiembre de
2026}

\begin{quick_sum}
    Antes de aplicar cualquier clasificador hay que preparar los
    datos: limpiarlos y ponerlos en una escala común, decidiendo entre
    normalizar y estandarizar. Luego, cómo partir la muestra para
    entrenar sin caer en overfitting ni en underfitting, y la variante
    estratificada de la validación cruzada para datos desbalanceados.
    Al final se abrió el tema nuevo, las redes neuronales artificiales,
    con la neurona biológica que les sirve de modelo.
\end{quick_sum}

\subsection*{Preparación de los datos}

A veces queremos aplicar un clasificador, un algoritmo que nos permita
separar, clasificar o estimar. Pero antes están los datos, y lo
primero es limpiarlos: eliminar columnas vacías, llenar los valores
faltantes y demás.

Después viene la escala. Los atributos suelen venir en unidades
distintas, y para evitar que unos dominen a otros hay que
\textbf{normalizar} o \textbf{estandarizar}.

\begin{center}
\begin{tabular}{lll}
  \toprule
  \textbf{Técnica} & \textbf{Qué hace} & \textbf{Fórmula} \\
  \midrule
  Estandarizar & Media $0$, varianza $1$
               & $z = \dfrac{x - \mu}{\sigma}$ \\
  \addlinespace
  Normalizar   & Todo entre $0$ y $1$
               & $x' = \dfrac{x - \min}{\max - \min}$ \\
  \bottomrule
\end{tabular}
\end{center}

Estandarizar reordena los datos en función de su media y su varianza.
Eso uniforma el manejo, para que ninguna variable se salga del rango
de las demás, y de paso aplasta a los datos atípicos: reduce esa
diferencia en vez de dejar que arrastre todo.

Escoger entre las dos es una de las decisiones que hay que tomar, y
depende de los datos y del objetivo. La regla que se dio en clase:

\begin{center}
\begin{tabular}{ll}
  \toprule
  \textbf{Distribución} & \textbf{Qué aplicar} \\
  \midrule
  No normal & Estandarizar \\
  Normal    & Normalizar \\
  \bottomrule
\end{tabular}
\end{center}

\begin{idea}
  La razón de fondo es de dónde salen las constantes de cada fórmula.
  Normalizar usa el mínimo y el máximo, que son los dos estadísticos
  más sensibles que existen a un valor extremo: si la distribución
  tiene cola larga, un solo punto lejano estira el denominador y
  comprime a todos los demás en una franja diminuta. Estandarizar usa
  media y desviación, que no se dejan mover tanto por un punto. Por
  eso la distribución no normal pide estandarizar. Si la distribución
  es normal, el mínimo y el máximo se portan bien y normalizar es
  seguro.
\end{idea}

\begin{ojo}
  El escalado se aplica \textbf{por variable}, no sobre todos los
  datos juntos. Cada columna tiene su propia media y desviación, o su
  propio mínimo y máximo. Mezclarlas destruye justo lo que se quería
  arreglar.
\end{ojo}

Todo esto sirve para que en el clasificador lineal no haya
descompensaciones en los pesos. Si un atributo va de $0$ a $1$ y otro
de $0$ a $100\,000$, el peso del segundo tiene que volverse diminuto
para compensar, y el entrenamiento se vuelve lento y mal condicionado.

\subsection*{Cómo dividir el conjunto de muestra}

La siguiente decisión es cómo partir la muestra para el proceso de
entrenamiento. Los dos extremos fallan por razones opuestas:

\begin{itemize}[itemsep=0.2em]
  \item Si le pasamos \textbf{todos} los datos al clasificador, se los
        va a aprender y no va a poder generalizar. Eso es
        \textbf{overfitting}, y además nos quedamos sin manera de
        medirlo, porque ya no hay datos que el modelo no haya visto.
  \item Si le pasamos \textbf{muy pocos}, no alcanza a aprender el
        patrón de manera adecuada. Eso es \textbf{underfitting}.
\end{itemize}

De ahí las tres particiones, cada una con su trabajo:

\begin{center}
\begin{tabular}{ll}
  \toprule
  \textbf{Partición} & \textbf{Para qué} \\
  \midrule
  Entrenamiento & Ajustar los parámetros del modelo \\
  Validación    & Escoger el modelo y sus hiperparámetros \\
  Prueba        & Estimar el error real, una sola vez \\
  \bottomrule
\end{tabular}
\end{center}

Los repartos usuales son $70/15/15$ o $80/10/10$ por ciento.

\begin{ojo}
  El conjunto de prueba se toca una vez, al final. Si se usa para
  decidir algo, deja de ser una estimación honesta del error: se
  convierte en parte del entrenamiento, y la generalización que
  reporta ya está contaminada.
\end{ojo}

\subsubsection*{Validación cruzada estratificada}

La otra opción es la validación cruzada, que ya se vio en la clase
del 3 de septiembre (figura \ref{fig:kfolds}): se parte el conjunto en
$k$ folds y se rota cuál se deja fuera.

Con datos desbalanceados eso no basta. Se parten \textbf{por clases} y
los folds se arman con datos de cada clase, de modo que cada fold
conserve la proporción original.

\begin{idea}
  Sin estratificar, con una clase rara y folds aleatorios puede tocar
  un fold sin un solo ejemplo de esa clase. Ese fold no mide nada
  sobre ella, y el promedio de los $k$ errores sale optimista por una
  razón que no tiene que ver con el clasificador.
\end{idea}

\subsection*{Redes neuronales artificiales}

Hasta aquí todos los modelos salieron de plantear un problema de
optimización. Pero alguien preguntó por qué no tratamos de emular el
modo en que aprende el cerebro humano, y de ahí nace la idea de las
redes neuronales artificiales.

\subsubsection*{La neurona biológica}

El cerebro está conectado por neuronas. Las conexiones entre neuronas
se llaman \textbf{sinapsis}, y se dan entre axones y dendritas.

\begin{center}
\begin{tabular}{ll}
  \toprule
  \textbf{Parte} & \textbf{Qué hace} \\
  \midrule
  Dendritas & Reciben las señales de otras neuronas \\
  Soma      & Integra lo que llega y decide si dispara \\
  Axón      & Lleva la señal de salida hacia afuera \\
  Sinapsis  & Une el axón de una con la dendrita de otra, y es
              donde la señal se pondera \\
  \bottomrule
\end{tabular}
\end{center}

Por qué es eficiente el cerebro? Por el número: del orden de $10^{11}$
neuronas y unas $10^{14}$ conexiones entre ellas. No es que cada
neurona sea potente, es que son muchísimas y están densamente
conectadas.

\subsubsection*{La neurona artificial}

El modelo artificial se queda con la estructura y tira la biología.
Cada parte tiene su contraparte:

\begin{center}
\begin{tabular}{ll}
  \toprule
  \textbf{Biológica} & \textbf{Artificial} \\
  \midrule
  Dendritas & Las entradas $x_1, \ldots, x_d$ \\
  Sinapsis  & Los pesos $w_1, \ldots, w_d$ \\
  Soma      & La suma ponderada más el sesgo \\
  Axón      & La salida, pasada por una función de activación \\
  \bottomrule
\end{tabular}
\end{center}

O sea
\[
    y = g\!\left( \sum_{i=1}^{d} w_i x_i + w_0 \right) ,
\]
con $g$ la función de activación.

\begin{idea}
  Con $g$ el signo, esa ecuación es exactamente el perceptrón de la
  clase del 27 de agosto. La neurona artificial no es un modelo nuevo,
  es el mismo clasificador lineal contado desde la biología. Lo nuevo
  viene de conectar muchas en capas, que es donde deja de ser lineal.
\end{idea}



\newpage

% =====================================================================
%  Clase - Neurona artificial y separabilidad lineal
% =====================================================================

\clase{Neurona artificial y separabilidad lineal}{15 de septiembre de
2026}

\begin{quick_sum}
    Continuación de la neurona. El TLU de McCulloch y Pitts con su
    ejemplo numérico, por qué el escalón es demasiado rígido para
    modelar una neurona real y cómo lo suaviza la sigmoide, el
    catálogo de funciones de activación y por qué importa que sean
    derivables. Al final, la lectura geométrica: el TLU es un
    clasificador lineal y su frontera es un hiperplano perpendicular
    al vector de pesos.
\end{quick_sum}

\subsection*{El TLU de McCulloch y Pitts}

La tarea es modelar algunos de los ingredientes de una red neuronal
biológica. El carácter de \emph{todo o nada} del potencial de acción
se caracteriza con una señal de dos valores, $0$ y $1$. Si un nodo
recibe $n$ señales de entrada $x_1, \ldots, x_n$, cada una solo puede
valer $0$ o $1$.

El efecto modulador de cada sinapsis se encapsula multiplicando la
señal que entra por un peso. Con $n$ pesos $w_1, \ldots, w_n$ se
forman los $n$ productos $w_i x_i$, y cada producto es el análogo de
un \textbf{potencial postsináptico}, que puede ser positivo o
negativo según el signo del peso.

Esos productos se combinan sumándolos, y eso emula lo que pasa en el
axón. El resultado es la \textbf{activación}:
\[
    a = \sum_{i=1}^{n} w_i x_i .
\]

Para emular la generación del potencial de acción hace falta un
\textbf{umbral} $\theta$: si la activación lo alcanza o lo supera, el
nodo emite $1$; si no, emite $0$.
\[
    y = \begin{cases}
      1 & \text{si } a \geq \theta , \\
      0 & \text{si } a < \theta .
    \end{cases}
\]

\begin{definicion}
  Esa unidad, suma ponderada más umbral con función escalón, se llama
  \textbf{TLU}, \emph{threshold logic unit}, y es de McCulloch y
  Pitts (1943).
\end{definicion}

\begin{ejemplo}
  Una unidad de cinco entradas con pesos
  $(0.5,\; 1.0,\; -1.0,\; -0.5,\; 1.2)$, a la que se le presenta el
  patrón $(1, 1, 1, 0, 0)$.

  La activación es
  \begin{align*}
    a &= (0.5 \times 1) + (1.0 \times 1) + (-1.0 \times 1)
       + (-0.5 \times 0) + (1.2 \times 0) \\
      &= 0.5 + 1.0 - 1.0 + 0 + 0 \\
      &= 0.5 .
  \end{align*}
  Con umbral $\theta = 0.2$, como $a = 0.5 > 0.2$, la salida es
  $y = 1$.

  Fíjate en el tercer término: el peso $w_3 = -1.0$ es negativo, así
  que esa entrada \textbf{inhibe} en vez de excitar, y se come
  completo el aporte de la segunda. Los pesos negativos son la
  contraparte de las sinapsis inhibitorias.
\end{ejemplo}

\subsection*{Comunicación con señales no binarias}

El escalón es demasiado rígido para lo que hace una neurona real. Se
cree que las neuronas codifican su señal en los \emph{patrones} de
disparo del potencial de acción, no en la simple presencia o ausencia
de un pulso. Si $f$ es la frecuencia de disparo, está acotada por
abajo en cero y por arriba en un $f_{\max}$ que impone la duración
del periodo refractario entre disparos.

La salida es suavizar el escalón hasta una función continua que
aplaste el rango, de modo que $y$ dependa de forma suave de $a$. La
forma conveniente es la \textbf{sigmoide logística}:
\[
    y = \sigma(a) = \frac{1}{1 + \exp\!\left(
        -\dfrac{a - \theta}{\rho} \right)} .
\]

Desglosando la fórmula pieza por pieza:

\begin{itemize}[itemsep=0.25em]
  \item El término $a - \theta$ centra la curva en el umbral: cuando
        $a = \theta$ el exponente es cero, $\exp(0) = 1$ y queda
        $y = 1/2$. El umbral es el punto donde la neurona está
        indecisa.
  \item El signo menos hace que la curva sea creciente. Si $a$ crece
        mucho, el exponente se va a $-\infty$, $\exp(-\infty) = 0$ y
        $y \to 1$. Si $a$ se hace muy negativa, el exponente se va a
        $+\infty$ y $y \to 0$.
  \item El divisor $\rho$ controla la forma. Valores grandes de
        $\rho$ achican el exponente y aplanan la curva; valores
        pequeños la hacen subir más abruptamente. En el límite
        $\rho \to 0$ se recupera el escalón. En muchos textos se
        omite, lo que equivale a fijar $\rho = 1$.
\end{itemize}

\begin{ojo}
  La sigmoide tiende a $1$ cuando $a$ crece, pero \textbf{nunca lo
  alcanza}. Igual por abajo: se acerca a $0$ sin tocarlo. Nunca
  satura del todo, y eso va a importar cuando se derive, porque la
  pendiente nunca es exactamente cero.
\end{ojo}

\begin{idea}
  La derivada de la sigmoide se escribe en términos de ella misma:
  \[
      \sigma'(a) = \sigma(a)\,\left[1 - \sigma(a)\right] .
  \]
  Eso no es una curiosidad, es lo que hace barato el entrenamiento:
  para calcular la pendiente en un punto no hay que evaluar nada
  nuevo, se reusa la salida que ya se calculó en el paso hacia
  adelante. Nota que vale cero en los extremos, donde $\sigma$ es $0$
  o $1$, y es máxima en $\sigma = 1/2$, o sea justo en el umbral.
\end{idea}

\subsection*{Elegir la función de activación}

Hay varias opciones, y cada una cambia el comportamiento de la
unidad.

\begin{center}
\begin{tabular}{lll}
  \toprule
  \textbf{Función} & \textbf{Definición} & \textbf{Rango} \\
  \midrule
  Escalón & $y = \begin{cases} 1 & a \geq \theta \\
                               0 & a < \theta \end{cases}$
          & $\{0, 1\}$ \\
  \addlinespace
  Sigmoide & $y = \dfrac{1}{1 + e^{-a}}$ & $(0, 1)$ \\
  \addlinespace
  ReLU & $y = \begin{cases} 0 & a < 0 \\
                            a & \text{si no} \end{cases}$
       & $[0, \infty)$ \\
  \addlinespace
  Tangente hiperbólica
       & $y = \dfrac{e^{a} - e^{-a}}{e^{a} + e^{-a}}$
       & $(-1, 1)$ \\
  \addlinespace
  Softmax & $Y_k = \dfrac{e^{a_k}}{\sum_i e^{a_i}}$
          & $(0,1)$, suman $1$ \\
  \bottomrule
\end{tabular}
\end{center}

Qué aporta cada una:

\begin{itemize}[itemsep=0.25em]
  \item \textbf{Escalón.} Decisión dura, sin grises. Es lo que define
        al TLU original.
  \item \textbf{Sigmoide.} El suavizado del escalón. Su salida en
        $(0,1)$ se puede leer como una probabilidad.
  \item \textbf{ReLU.} Deja pasar lo positivo y corta lo negativo. Es
        barata de calcular y su derivada es $1$ o $0$, sin aplastarse
        para valores grandes.
  \item \textbf{Tangente hiperbólica.} Es una sigmoide centrada en
        cero, con salida en $(-1,1)$. Que esté centrada ayuda porque
        las salidas no quedan todas del mismo signo.
  \item \textbf{Softmax.} No opera sobre una neurona sino sobre toda
        la capa: normaliza las activaciones para que sumen $1$. Es la
        que se usa en la capa de salida de un clasificador
        multiclase.
\end{itemize}

\begin{ojo}
  El punto que hay que retener: \textbf{el escalón no es derivable}.
  Las otras sí, y eso es esencial para el cálculo eficiente de los
  pesos. Todo el entrenamiento por descenso de gradiente descansa en
  poder derivar la función de activación, así que el escalón deja a
  la unidad sin manera práctica de aprender.
\end{ojo}

% =====================================================================

\subsection*{El TLU como clasificador lineal}

Un TLU parte sus patrones de entrada en dos categorías según su
respuesta binaria. Esas categorías se pueden ver como regiones de un
espacio multidimensional separadas por el equivalente en muchas
dimensiones de una recta o un plano.

\begin{ejemplo}
  Un TLU de dos entradas con $w_1 = 1$, $w_2 = 1$ y $\theta = 1.5$.
  Las cuatro entradas booleanas posibles dan:

  \begin{center}
  \begin{tabular}{cccc}
    \toprule
    $x_1$ & $x_2$ & \textbf{Activación} & \textbf{Salida} \\
    \midrule
    $0$ & $0$ & $0$ & $0$ \\
    $0$ & $1$ & $1$ & $0$ \\
    $1$ & $0$ & $1$ & $0$ \\
    $1$ & $1$ & $2$ & $1$ \\
    \bottomrule
  \end{tabular}
  \end{center}

  Solo el patrón $(1,1)$ supera el umbral de $1.5$. El TLU está
  clasificando sus entradas en dos grupos: las que dan $1$ y las que
  dan $0$. Como cada patrón tiene dos componentes, se puede dibujar
  cada uno como un punto del \textbf{espacio de patrones}, y ahí
  $(1,1)$ queda de un lado de una recta y los otros tres del otro.

  Esa función es el AND lógico.
\end{ejemplo}

\subsubsection*{De dónde sale la recta}

La condición crítica de clasificación ocurre cuando la activación
iguala al umbral. Para dos entradas, igualando $a$ y $\theta$:
\begin{align*}
  w_1 x_1 + w_2 x_2 &= \theta \\
  w_2 x_2           &= -w_1 x_1 + \theta \\
  x_2               &= -\left(\frac{w_1}{w_2}\right) x_1
                       + \left(\frac{\theta}{w_2}\right) .
\end{align*}

Eso es una recta $x_2 = a x_1 + b$, con pendiente
$a = -w_1/w_2$ y ordenada al origen $b = \theta/w_2$ sobre el eje
$x_2$. Con los valores del ejemplo, $w_1 = w_2 = 1$ y
$\theta = 1.5$, queda $x_2 = -x_1 + 1.5$: una recta que corta ambos
ejes en $1.5$ y deja a $(1,1)$ arriba, solo.

\begin{idea}
  Los pesos fijan la \textbf{inclinación} de la frontera y el umbral
  fija su \textbf{posición}. Entrenar es mover esa recta hasta que
  los puntos queden del lado correcto. Por eso a los TLU se les llama
  \textbf{clasificadores lineales} y a sus patrones,
  \textbf{linealmente separables}. Es el mismo objeto que el
  perceptrón de la clase del 27 de agosto, visto desde la biología.
\end{idea}

\subsubsection*{En forma vectorial}

Usando la definición de producto interno, la activación de un TLU de
$n$ entradas se escribe
\[
    a = W \cdot X ,
\]
y la condición crítica queda
\[
    W \cdot X = \theta .
\]

Esa ecuación define un \textbf{hiperplano perpendicular al vector de
pesos}. De un lado quedan todos los patrones que el TLU clasifica
como $1$; del otro, los que clasifica como $0$.

\begin{ojo}
  Que el hiperplano sea perpendicular a $W$ no es un detalle
  decorativo: quiere decir que $W$ apunta hacia el lado de los
  positivos. Esa observación es la que va a justificar toda la regla
  de entrenamiento de la clase siguiente, porque corregir un error
  será \emph{rotar} $W$ hacia el patrón o alejarlo de él.
\end{ojo}

\subsection*{El umbral como un peso más}

Hay un truco de notación que simplifica todo lo que viene. Partiendo
de la condición de salida $W \cdot X \geq \theta$, se resta $\theta$
de ambos lados y se hace explícito el signo:
\[
    W \cdot X + (-1)\,\theta \geq 0 .
\]

Leído así, el umbral es \textbf{un peso más}, alimentado por una
entrada fija en $-1$. El vector de pesos, que era de dimensión $n$
para una unidad de $n$ entradas, pasa a ser el vector de dimensión
$n+1$
\[
    [\,w_1,\; w_2,\; \ldots,\; w_n,\; \theta\,] ,
\]
que se llama \textbf{vector de pesos aumentado}. Con él, la función
del nodo se escribe sin umbral a la vista:
\[
    W \cdot X \geq 0 \;\Longrightarrow\; y = 1 ,
    \qquad
    W \cdot X < 0 \;\Longrightarrow\; y = 0 .
\]

\begin{idea}
  En el espacio aumentado, $W \cdot X = 0$ define el hiperplano de
  decisión, que sigue siendo ortogonal a $W$. Y la condición de
  umbral cero significa que el hiperplano \textbf{pasa por el
  origen}, porque esa es la única forma de que $W \cdot X = 0$ tenga
  solución en el origen. La ganancia práctica es que ya no hay que
  entrenar el umbral aparte: entra al mismo saco que los pesos y la
  misma regla lo ajusta.
\end{idea}

\newpage

% =====================================================================
%  Clase - Entrenamiento, regla delta y backpropagation
% =====================================================================

\clase{Descenso de gradiente y backpropagation}{17 de septiembre de
2026}

\begin{quick_sum}
    Cómo se entrena. Primero la regla del perceptrón, que sale de un
    argumento puramente geométrico sobre rotar el vector de pesos.
    Luego el descenso de gradiente sobre una función de error, la
    regla delta con su derivación, y por qué converge donde el
    perceptrón oscila. Al final el problema del XOR, que obliga a
    meter capas ocultas, y backpropagation como la respuesta a quién
    tiene la culpa del error.
\end{quick_sum}

\subsection*{Entrenar una red}

Ajustar los pesos y los umbrales se hace con un proceso iterativo: se
presentan ejemplos de la tarea una y otra vez, y en cada presentación
se hacen cambios pequeños para acercarlos a sus valores deseados. A
eso se le llama \textbf{entrenar la red}, y al conjunto de ejemplos,
\textbf{conjunto de entrenamiento}.

Trabajando con un solo nodo, el conjunto de entrenamiento son pares
$\{V, t\}$, donde $V$ es un vector de entrada y $t$ es la clase
objetivo, $1$ o $0$, a la que $V$ pertenece. La condición de salida
$1$ se escribe $W \cdot V \geq \theta$.

\subsection*{Ajustar el vector de pesos}

Aquí el razonamiento es geométrico y vale la pena seguirlo despacio,
porque de él sale la regla.

\textbf{Caso 1: el objetivo es $t = 1$ pero sale $y = 0$.} Para
producir un $0$ la activación tuvo que ser negativa, así que el
producto punto $W \cdot V$ fue negativo. Eso significa que los dos
vectores apuntaban en direcciones opuestas, con ángulo mayor a
$90^{\circ}$. Lo que hace falta es \textbf{rotar $W$ para que apunte
más hacia $V$}, y eso se logra sumándole una fracción de $V$:
\[
    W' = W + \alpha V , \qquad 0 < \alpha < 1 .
\]

\textbf{Caso 2: el objetivo es $t = 0$ pero sale $y = 1$.} Ahora la
activación fue positiva cuando debió ser negativa, el ángulo es menor
a $90^{\circ}$, y hay que \textbf{alejar $W$ de $V$} restando una
fracción:
\[
    W' = W - \alpha V .
\]

Las dos situaciones se combinan en una sola regla:
\[
    \boxed{\;W' = W + \alpha\,(t - y)\,V\;}
\]

\begin{idea}
  El factor $(t - y)$ hace todo el trabajo de selección, y solo puede
  tomar tres valores. Si acertó, $t - y = 0$ y no se toca nada. Si
  $t=1$ y $y=0$, vale $+1$ y se suma. Si $t=0$ y $y=1$, vale $-1$ y
  se resta. Una sola expresión cubre los tres casos sin condicionales.
\end{idea}

Esa es la \textbf{regla de entrenamiento del perceptrón}, de Frank
Rosenblatt (1957-1958). El $\alpha$ se llama \textbf{tasa de
aprendizaje} y controla qué tan grande es cada corrección.

\begin{algorithm}[H]
\caption{Aprendizaje del perceptrón}
\begin{algorithmic}[1]
\Repeat
  \For{cada par de entrenamiento $(V, t)$}
    \State evaluar la salida $y$ al presentar $V$ al TLU
    \If{$y \neq t$}
      \State $W \gets W + \alpha\,(t - y)\,V$
    \Else
      \State no hacer nada
    \EndIf
  \EndFor
\Until{$y = t$ para todos los vectores}
\end{algorithmic}
\end{algorithm}

% =====================================================================

\subsection*{Descenso de gradiente sobre un error}

La regla anterior sale de la geometría. Ahora se llega a una parecida
desde la optimización, y esa segunda ruta es la que generaliza a
redes de varias capas.

Considerando una red de un solo TLU, cualquier función de error $E$
que exprese la discrepancia entre la salida deseada y la real puede
verse como una \textbf{función de los pesos}:
\[
    E = E(w_1, w_2, \ldots, w_{n+1}) .
\]
El $n+1$ es porque el umbral ya va dentro, como peso aumentado.

El vector de pesos óptimo se encuentra minimizando $E$ por descenso
de gradiente:
\[
    \Delta w_i = -\,\alpha\,\frac{\partial E}{\partial w_i} .
\]

\begin{idea}
  El signo menos es el corazón del método. La derivada
  $\partial E / \partial w_i$ dice hacia dónde el error
  \textbf{crece} si muevo ese peso. Como queremos bajarlo, nos
  movemos en la dirección contraria. Si la pendiente es positiva, el
  peso baja; si es negativa, sube. Y donde la pendiente es cero, en
  el mínimo, el peso deja de moverse solo.
\end{idea}

Si $e^p$ es el error del patrón de entrenamiento $p$, el error total
es el promedio sobre todos los patrones:
\[
    E = \frac{1}{N} \sum_{p=1}^{N} e^p .
\]

\subsubsection*{Cómo definir el error de un patrón}

Vale la pena ver los tres intentos, porque cada uno arregla un
problema del anterior.

\textbf{Primer intento: la diferencia.} $e^p = t^p - y^p$. No sirve:
la combinación $t^p = 0$, $y^p = 1$ da $-1$, mientras
$t^p = 1$, $y^p = 0$ da $+1$. Los dos casos están igual de
equivocados, pero uno cuenta como error negativo y al promediarlos se
cancelarían.

\textbf{Segundo intento: elevar al cuadrado.}
$e^p = (t^p - y^p)^2$. Ya es simétrico y siempre positivo. Pero hay
un problema más fino: $y^p$ depende de $a^p$ a través de la
\textbf{función escalón, que es discontinua}, así que no se puede
derivar.

\textbf{Tercer intento: medir contra la activación.} Como la salida
cambia justo cuando la activación cambia de signo, con $a \geq 0$
dando $y = 1$, se define el error respecto a la activación:
\[
    e^p = (t^p - a^p)^2 .
\]
Y se agrega un factor de $1/2$, que no cambia dónde está el mínimo
pero simplifica las derivadas:
\[
    e^p = \tfrac{1}{2}\,(t^p - a^p)^2 .
\]

% =====================================================================

\subsection*{La regla delta}

Como $E$ depende de todos los patrones, sus derivadas también, y eso
obligaría a presentar el conjunto completo para evaluar un solo
gradiente. Se evita presentando un patrón $p$, evaluando
$\partial e^p / \partial w_i$ y usándolo como \textbf{estimación} del
gradiente verdadero $\partial E / \partial w_i$.

\subsubsection*{La derivada, paso a paso}

Las diapositivas dicen que se puede demostrar. Se demuestra así.
Partiendo de $e^p = \frac{1}{2}(t^p - a^p)^2$ con
$a^p = \sum_i w_i x_i^p$, y aplicando la regla de la cadena:
\begin{align*}
  \frac{\partial e^p}{\partial w_i}
    &= \frac{1}{2} \cdot 2\,(t^p - a^p) \cdot
       \frac{\partial (t^p - a^p)}{\partial w_i}
    && \text{cadena sobre el cuadrado} \\
    &= (t^p - a^p) \cdot
       \left( - \frac{\partial a^p}{\partial w_i} \right)
    && t^p \text{ es constante} \\
    &= -\,(t^p - a^p)\,x_i^p
    && \frac{\partial}{\partial w_i}\sum_j w_j x_j^p = x_i^p .
\end{align*}

Ahí se ve para qué era el $1/2$: se cancela con el $2$ que baja del
exponente y la expresión queda limpia.

Sustituyendo en $\Delta w_i = -\alpha\,\partial E/\partial w_i$, los
dos signos menos se cancelan y queda la \textbf{regla delta}:
\[
    \boxed{\;\Delta w_i = \alpha\,(t^p - a^p)\,x_i^p\;}
\]

\subsubsection*{Cómo se lee la fórmula}

El gradiente depende de dos cosas, y las dos tienen sentido:

\begin{itemize}[itemsep=0.25em]
  \item \textbf{De la discrepancia $(t^p - a^p)$.} Entre más grande
        sea el error, más grande la corrección, y viceversa. Si ya
        acertó, la corrección es cero.
  \item \textbf{De la entrada $x_i^p$.} Si esa entrada vale cero, la
        $i$-ésima componente no está contribuyendo a la activación, y
        entonces su peso no tiene culpa de nada: no se toca. Si la
        entrada es grande, ese peso sí influye mucho y se corrige
        proporcionalmente.
\end{itemize}

\begin{algorithm}[H]
\caption{Entrenamiento con la regla delta}
\begin{algorithmic}[1]
\Repeat
  \For{cada par de entrenamiento $(V, t)$}
    \State evaluar la activación $a$ al presentar $V$ al TLU
    \State $\Delta w_i \gets \alpha\,(t^p - a^p)\,x_i^p$
           \Comment{para cada peso $i$}
    \State $w_i \gets w_i + \Delta w_i$
  \EndFor
\Until{la razón de cambio del error sea suficientemente pequeña}
\end{algorithmic}
\end{algorithm}

Esta regla la propusieron Widrow y Hoff (1960), así que originalmente
se llamó \textbf{regla de Widrow-Hoff}. Se conoce más como
\textbf{regla delta}, y al término $(t^p - a^p)$ se le dice
\emph{la} $\delta$, porque es una diferencia.

\begin{ojo}
  La diferencia entre las dos reglas. Si \textbf{no existe solución},
  o sea si el problema no es linealmente separable, bajo la regla del
  perceptrón los pesos se siguen alterando en cantidades
  significativas y el vector \textbf{oscila} sin parar. La regla
  delta, en cambio, siempre converge (con $\alpha$ suficientemente
  chica) a un vector $w_0$ donde el error es mínimo, aunque ese
  mínimo no sea cero y algunos patrones queden mal clasificados.
  Converger a la mejor respuesta imperfecta es mucho mejor que no
  converger.
\end{ojo}

\subsubsection*{Versión con sigmoide}

Se puede reinstalar la salida $y^p$ en la definición del error,
ahora que se usa una sigmoide continua en vez del escalón:
\[
    e^p = \tfrac{1}{2}\,(t^p - y^p)^2 ,
    \qquad y^p = \sigma(a^p) .
\]
Al derivar aparece un eslabón más en la cadena, la derivada de la
sigmoide:
\[
    \frac{\partial e^p}{\partial w_i}
    = -\,(t^p - y^p) \cdot
      \underbrace{\frac{d\sigma}{da}}_{\sigma'(a)} \cdot\,
      x_i^p ,
\]
y por lo tanto
\[
    \boxed{\;\Delta w_i
    = \alpha\,\sigma'(a)\,(t^p - y^p)\,x_i^p\;}
\]

\begin{idea}
  Ese factor $\sigma'(a)$ gobierna todo. El efecto de cambiar un peso
  depende de qué tan sensible sea la salida a la activación, y esa
  sensibilidad es justo la pendiente de la sigmoide. En las colas,
  donde la sigmoide está casi plana, $\sigma'(a) \approx 0$ y el peso
  casi no se mueve aunque el error sea grande: la neurona está
  saturada. El aprendizaje es más rápido cerca del umbral, donde la
  pendiente es máxima.
\end{idea}

% =====================================================================

\subsection*{El límite: el XOR}

Todo lo anterior descansa en que los patrones sean linealmente
separables. Hay funciones booleanas elementales que no lo son.

\begin{ejemplo}
  Tomando $(0,0)$ y $(1,1)$ con $t = 0$, y $(0,1)$ y $(1,0)$ con
  $t = 1$. Los positivos quedan en dos esquinas opuestas del cuadrado
  y los negativos en las otras dos.

  \textbf{No existe ninguna recta} que separe unos de otros: hay que
  dejar dos esquinas diagonalmente opuestas de un lado, y cualquier
  recta que deje a $(0,1)$ y $(1,0)$ juntas necesariamente encierra
  también a una de las otras dos.

  Esa función es el XOR, y el señalamiento es de Minsky y Papert en
  \emph{Perceptrons} (1969).
\end{ejemplo}

\begin{ojo}
  Ese libro frenó la investigación en redes neuronales casi dos
  décadas. La limitación que señalaba era real, pero era la de
  \emph{una sola} unidad, no la de la idea entera.
\end{ojo}

\subsection*{La salida: más capas}

La solución se encuentra fácil si el espacio se parte en
\textbf{tres regiones} en vez de dos. Con dos rectas paralelas se
puede dejar a $(0,1)$ y $(1,0)$ en la franja de en medio y a las
otras dos esquinas afuera, una de cada lado.

\begin{enumerate}[itemsep=0.25em]
  \item Si en lugar de una neurona de salida se usan \textbf{dos},
        se obtienen dos rectas, y con ellas quedan delimitadas las
        tres regiones.
  \item Para escoger entre una zona y otra de las tres hace falta
        \textbf{otra capa}, con una neurona cuyas entradas sean las
        salidas de las neuronas anteriores.
\end{enumerate}

Eso da una red de tres neuronas repartidas en dos capas: dos
neuronas que reciben $x_1$ y $x_2$ con pesos $w_{11}, w_{21}, w_{12},
w_{22}$, y una tercera que recibe las salidas de esas dos con pesos
$w_{31}, w_{32}$ y produce $y$.

\begin{definicion}
  Una capa está \textbf{totalmente conectada} cuando las entradas de
  cada unidad de la capa son las mismas. Una capa toma una entrada
  $x \in \mathbb{R}^m$ y produce una salida $y \in \mathbb{R}^n$.
\end{definicion}

Con $m$ neuronas en una capa, el error se suma sobre todas:
\[
    e^p = \frac{1}{2} \sum_{j=1}^{m} (t_j^p - y_j^p)^2 ,
\]
y el gradiente del $i$-ésimo peso de la $j$-ésima neurona solo
involucra el término de error de \emph{esa} neurona:
\[
    \Delta w_{ji} = \alpha\,\sigma'(a_j)\,(t_j^p - y_j^p)\,x_{ji}^p .
\]

% =====================================================================

\subsection*{Backpropagation}

La idea sigue siendo descenso de gradiente sobre $E$ como función de
los pesos, pero ahora hay que tomar en cuenta los pesos de dos capas.

Para la \textbf{capa de salida} es directo: se tiene acceso al
desajuste entre el objetivo y la salida, y la fórmula de arriba
aplica tal cual.

El problema son las neuronas de la \textbf{capa oculta}, y la
pregunta se plantea así: cuánta \emph{culpa} o responsabilidad hay
que asignarle a los nodos ocultos por el error que aparece en la
salida. Porque claramente, si la capa oculta le está pasando mala
información a la de salida, las neuronas de salida no pueden aspirar
a acertar.

\subsubsection*{Descomponer la regla conocida}

El punto de partida es mirar la regla que ya se tiene y separarla en
piezas:
\[
    \Delta w_{ji} = \alpha\;
    \underbrace{\sigma'(a_j)}_{\text{pendiente}}\;
    \underbrace{(t_j^p - y_j^p)}_{\delta^j}\;
    \underbrace{x_{ji}^p}_{\text{entrada}} .
\]

La entrada y la pendiente de la sigmoide están determinadas
únicamente por la \textbf{estructura de la neurona}, así que van a
aparecer igual en la expresión de una neurona oculta. El término que
sobra, $\delta = (t^p - y^p)$, es el específico de las neuronas de
salida, y toda la tarea se reduce a encontrar su equivalente para una
neurona oculta.

Con eso, la regla de la $k$-ésima neurona oculta se escribe ya:
\[
    \Delta w_{ki} = \alpha\,\sigma'(a_k)\,\delta^k\,x_{ki}^p ,
\]
a falta de saber qué es $\delta^k$.

\subsubsection*{Encontrar la delta de una neurona oculta}

Considera el enlace entre la $k$-ésima neurona oculta y la
$j$-ésima de salida. Hay dos factores en juego:

\begin{itemize}[itemsep=0.25em]
  \item La contribución que la neurona $j$ hace al error está
        expresada en su propia $\delta^j$, que ya se sabe calcular.
  \item La influencia que $k$ tiene sobre $j$ está dada por el peso
        $w_{jk}$ de la conexión entre ellas.
\end{itemize}

La interacción entre los dos se captura combinándolos
\textbf{multiplicativamente}, como $\delta^j w_{jk}$. Y como la
neurona $k$ casi con seguridad alimenta a varias neuronas de salida,
hay que sumar esos productos sobre todas:
\[
    \boxed{\;\delta^k = \sum_{j \in I_k} \delta^j\,w_{jk}\;}
\]
donde $I_k$ es el conjunto de neuronas de la capa de salida que
reciben entrada de la neurona oculta $k$.

\begin{idea}
  Ahí está el nombre del método. La $\delta$ de una neurona oculta no
  se mide contra ningún objetivo, porque no hay objetivo para una
  capa oculta. Se \textbf{arma hacia atrás}, juntando las deltas de
  las neuronas a las que alimenta, cada una pesada por la fuerza de
  la conexión. El error viaja de regreso por las mismas aristas que
  la señal usó para ir hacia adelante, y los pesos hacen de factor de
  reparto en las dos direcciones.
\end{idea}

\subsubsection*{La regla unificada}

Para cualquier neurona $k$, oculta o de salida, el descenso de
gradiente queda en una sola línea:
\[
    \Delta w_{ki} = \alpha\,\delta^k\,x_{ki}^p ,
\]
y lo único que cambia es cómo se calcula $\delta^k$:
\begin{align*}
  \delta^k &= \sigma'(a_k)\,(t_k^p - y_k^p)
    && \text{si } k \text{ es de salida,} \\
  \delta^k &= \sigma'(a_k) \sum_{j \in I_k} \delta^j\,w_{jk}
    && \text{si } k \text{ es oculta.}
\end{align*}

\begin{algorithm}[H]
\caption{Backpropagation}
\begin{algorithmic}[1]
\State presentar el patrón en la capa de entrada
\State que las neuronas ocultas evalúen su salida con ese patrón
\State que las neuronas de salida evalúen la suya con el resultado
       del paso anterior
\State aplicar el patrón objetivo a la capa de salida
\State calcular las $\delta$ de las neuronas de salida
\State entrenar cada neurona de salida por descenso de gradiente
\State para cada neurona oculta, calcular su $\delta^k$
\State entrenar cada neurona oculta con la $\delta^k$ del paso
       anterior, por descenso de gradiente
\end{algorithmic}
\end{algorithm}

\begin{center}
\begin{tabular}{lll}
  \toprule
  \textbf{Pasos} & \textbf{Nombre} & \textbf{Qué fluye} \\
  \midrule
  1 a 3 & Paso hacia adelante & La información, de entrada a salida \\
  4 a 8 & Paso hacia atrás    & Las $\delta$, de salida a ocultas \\
  \bottomrule
\end{tabular}
\end{center}

El paso 7 es el que propaga las $\delta$ de vuelta desde los nodos de
salida hacia las unidades ocultas, y de ahí el nombre
\textbf{backpropagation}. En los años ochenta lo popularizaron
Geoffrey Hinton, David Rumelhart y Ronald Williams.

\begin{idea}
  Se puede tener más de una capa oculta, y se entrena exactamente
  igual: las $\delta$ de cada capa, salvo la de salida, se expresan
  en términos de las de la capa que tienen encima. La recursión es la
  misma sin importar la profundidad, y eso es lo que vuelve
  entrenable a una red arbitrariamente honda.
\end{idea}


\end{document}

%############################# END OF LIBRETA.TEX ##############################
%###############################################################################
